\chapter{Diskrete-Fourier-Transformation}\label{sec:diskrete-fourier-transformation}

In Kapitel \hyperref[spektrum-von-signalfolgen]{7} wird die Fourier-Transformation von Signalfolgen eingeführt. Mit ihr können die Spektren von Signalen und die Frequenzgänge von Systemen berechnet werden. In der praktischen Anwendung zeigen sich allerdings zwei Schwächen:

\begin{itemize}
\item
  Die Fourier-Transformation von Signalfolgen eignet sich nur für analytisch definierte Signale, für die unendlich viele Folgenwerte vorliegen. Messergebnisse sind jedoch immer auf ein endliches Zeitintervall begrenzt, sodass das Spektrum des zugrunde liegenden Signals nur über den vorliegenden Signalausschnitt geschätzt werden kann.
\item
  Außerdem ist das Ergebnis der Fourier-Transformation von Signalfolgen ein kontinuierliches Spektrum \(X(Ω)\), das numerisch nicht oder nur aufwendig bearbeitet werden kann. Für die digitale Signalverarbeitung im Spektralbereich wäre ein diskretes Spektrum vorteilhaft, das wie bei der Fourier-Reihe aus einzelnen Koeffizienten besteht.
\end{itemize}

In diesem Kapitel wird die Diskrete-Fourier-Transformation (DFT) eingeführt. Sie erlaubt die Schätzung von Spektren auf Basis von Signalausschnitten. Das Ergebnis der Diskrete-Fourier-Transformation ist ein diskretes Spektrum. Nach ihrer Definition wird ihre Berechnung an einführenden Beispielen erläutert.

Die Eigenschaften der Diskreten-Fourier-Transformation entsprechen weitgehend denen der Fourier-Transformation für Signalfolgen. Abweichungen dazu werden herausgearbeitet und an Beispielen veranschaulicht.

Bei der Schätzung von Spektren analoger Signale muss die gesamte Signalkette von der Wandlung bis zur Berechnung der Fourier-Transformierten analysiert werden. Dabei wird auf den sogenannten Leakage-Effekt sowie die Auswirkung unterschiedlicher Fensterfunktionen eingegangen.

Die Auflösung der Diskreten-Fourier-Transformation ist abhängig von der Anzahl der Messwerte der Signalfolge. Die Ergänzung der Messwerte mit Nullen (Zero-Padding) im Zeitbereich entspricht einer Interpolation des Spektrums im Frequenzbereich. Der mathematische Hintergrund dazu wird erläutert und die Grenzen des Verfahrens werden an Beispielen diskutiert. Auf ähnliche Weise kann durch das Einfügen von Nullen im Frequenzbereich eine Interpolation im Zeitbereich erreicht werden.

Für den praktischen Einsatz der Diskreten-Fourier-Transformation wird sie als Fast-Fourier-Transformation implementiert. Das Ergebnis der beiden Transformationen ist unter der Annahme einer Signalfolge mit einer Länge von N = \(2^{n}\) identisch. Durch Ausnutzen von Periodizität und Symmetrien liegt das Ergebnis der Fast-Fourier-Transformation jedoch deutlich schneller vor als bei der Diskreten-Fourier-Transformation.

Die Fourier-Transformation erlaubt eine effiziente Implementierung der Faltung im Zeitbereich durch eine Multiplikation der diskreten Fourier-Transformierten im Frequenzbereich. Um Signale kontinuierlich zu bearbeiten, werden die Signale blockweise verarbeitet. Dazu kann zum Beispiel das Overlap-Add-Verfahren eingesetzt werden.

\section{Grundlagen der Diskreten-Fourier-Transformation}\label{sec:10.1}

\subsection{Definition der Diskreten-Fourier-Transformation}\label{sec:10.1.1}

Die Fourier-Koeffizienten \(A_{n}\) der Fourier-Reihe eines in \(T_{B}\) periodischen Signals \(x(t)\) ergeben sich aus
\begin{align}
A_{n} &= \frac{1}{T_{B}}  \int_{- T_{B}/2}^{T_{B}/2}{x(t)  e^{- j  2  \pi  \frac{n  t}{T_{B}}}\ dt} = \frac{1}{T_{B}}  \int_{0}^{T_{B}}{x(t)  e^{- j  2  \pi  \frac{n  t}{T_{B}}}\ dt} \label{eq:10.1}
\end{align}
Mit der Diskreten-Fourier-Transformation wird das Spektrum abgetasteter Signale beschrieben. Mit den Ausführungen in Kapitel \hyperref[signalabtastung-und-rekonstruktion]{2} kann ein ideal abgetastete Signal \(x_{A}(t)\) dargestellt werden als
\begin{align}
x_{A}(t) &= \sum\limits_{k = - \infty}^{\infty}{x[ k]  \delta\left( t - k  T_{A} \right)} \label{eq:10.2}
\end{align}
Die Signalfolge \(x[k]\) wird für einen Zeitraum \(T_{B}\) beobachtet.
\begin{align}
T_{B} &= N  T_{A} \label{eq:10.3}
\end{align}
Die zeitlich begrenzte Beobachtung wird mit einer Fensterfunktion \(w[k]\) beschrieben.
\begin{align}
x_{w}[ k] &= x[ k]  w[ k] = x[ k]  \left( \sigma[ k] - \sigma[ k - N] \right) \label{eq:10.4}
\end{align}
Zur Herleitung der Diskreten-Fourier-Transformation wird dieser Signalausschnitt periodisch in der Beobachtungszeit \(T_{B}\) wiederholt. Das periodische Signal hat die Fourier-Reihe
\begin{align}
A_{n} &= \frac{1}{T_{B}}  \int_{0}^{T_{B}}{x_{A}(t)  e^{- j  2  \pi  \frac{n  t}{T_{B}}}\ dt} = \frac{1}{T_{B}}  \int_{0}^{T_{B}}{\sum\limits_{k = 0}^{N - 1}{x_{W}[ k]  \delta\left( t - k  T_{A} \right)}  e^{- j  2  \pi  \frac{n  t}{T_{B}}}\ dt} \label{eq:10.5}
\end{align}
mit den Indizes
\begin{align}
n &= 0\ ...\ N - 1 \label{eq:10.6}
\end{align}
Vertauschen von Summation und Integration sowie Anwendung der Ausblendeigenschaft der Impulsfunktion führt zur Fourier-Reihe der periodisch wiederholten Signalfolge \(x_{W}\){[}k{]}
\begin{align}
{A_{n} } &=  {\frac{1}{T_{B} }  \sum\limits _{k=0}^{N-1}x_{W} \left[k\right] \int _{0}^{T_{B} }\delta \left(t-k T_{A} \right) e^{-j 2 \pi  \frac{n t}{T_{B} } } {\; }dt  =\frac{1}{T_{B} }  \sum\limits _{k=0}^{N-1}x_{W} \left[k\right] \int _{0}^{T_{B} }\delta \left(t-k T_{A} \right) e^{-j 2 \pi  \frac{n k T_{A} }{T_{B} } } {\; }dt  } \\
{} &=  {\frac{1}{N T_{A} }  \sum\limits _{k=0}^{N-1}x_{W} \left[k\right] e^{-j 2 \pi  \frac{n k T_{A} }{N T_{A} } }  \int _{0}^{T_{B} }\delta \left(t-k T_{A} \right){\; }dt  =\frac{1}{N T_{A} }  \sum\limits _{k=0}^{N-1}x_{W} \left[k\right] e^{-j 2 \pi  \frac{n k}{N} }  }
 \label{eq:10.7}
\end{align}
In Anlehnung an die Definitionsgleichung der Fourier-Transformation von Signalfolgen
\begin{align}
X_{W}(\Omega) &= \sum\limits_{k = - \infty}^{\infty}{x_{W}[ k]  e^{- j  \Omega  k}} \label{eq:10.8}
\end{align}
wird die diskrete Fourier-Transformation definiert als
\begin{align}
X_{n} &= \sum\limits_{k = 0}^{N - 1}{x_{W}[ k]  e^{- j  \left( 2  \pi  \frac{n}{N}  k \right)}}\ \label{eq:10.9}
\end{align}
beziehungsweise kurz
\begin{align}
X_{n} &= \sum\limits_{k = 0}^{N - 1}{x[ k]  e^{- j  \left( 2  \pi  \frac{n}{N}  k \right)}}\ \label{eq:10.10}
\end{align}
Aus einem Vergleich von den Gleichungen \eqref{eq:10.8} und \eqref{eq:10.9} wird deutlich, dass die Diskrete-Fourier-Transformation an den Frequenzen
\begin{align}
\Omega_{n} &= \frac{2  \pi  n}{N} \label{eq:10.11}
\end{align}
den Abtastwerten des Spektrums \(X_{W}\)(\(Ω_{n}\)) entspricht. Wegen des diskreten Spektrums wird die Transformation als Diskrete-Fourier-Transformation (DFT) bezeichnet. Sie geht von einer endlichen Anzahl von N Abtastwerten eines Signals aus und ist als endliche Summe definiert. Damit lässt sie sich in technischen Systemen verarbeiten.

Über die Umrechnung der Definitionsgleichung ergibt sich
\begin{align}
X_{n + N} &= \sum\limits_{k = 0}^{N - 1}{x[ k]  e^{- j  \left( 2  \pi  \frac{n + N}{N}  k \right)}} = \sum\limits_{k = 0}^{N - 1}{x[ k]  e^{- j  \left( 2  \pi  \frac{n}{N}  k \right)}  e^{- j  \left( 2  \pi  \frac{N}{N}  k \right)}} = \sum\limits_{k = 0}^{N - 1}{x[ k]  e^{- j  \left( 2  \pi  \frac{n}{N}  k \right)}} = X_{n}\ \label{eq:10.12}
\end{align}
Die Diskrete-Fourier-Transformierte ist periodisch in N.

\begin{example}{Rechteckfolge}
\label{beispiel-rechteckfolge}

Das Vorgehen bei der Diskreten-Fourier-Transformierten wird am Beispiel einer Rechteckfolge veranschaulicht. Sie besteht aus acht von null verschiedenen Abtastwerten und wird im Zeitraum von\\
0 \ldots{} 15·\(T_{A}\) beobachtet. Außerhalb dieses Zeitraums ist die Folge null. Dadurch entspricht das gefensterte Signal \(x_{W}\){[}k{]} dem Signal \(x[k]\).
\begin{align}
x[ k] &= \sigma[ k - 2] - \sigma[ k - 7] \label{eq:10.13}
\end{align}
Die Frequenzen, an denen die Diskrete-Fourier-Transformierte berechnet wird, ergeben sich mit dem Index \(n = 0\) \ldots{} 15 zu
\begin{align}
\omega_{n} &= \frac{\Omega_{n}}{T_{A}} = \frac{2  \pi  n}{T_{A}  N} \label{eq:10.14}
\end{align}
Die Diskrete-Fourier-Transformierte errechnet sich an diesen Stellen zu
\begin{align}
X_{n} &= \sum\limits_{k = 0}^{15}{x[ k]  e^{- j  \left( 2  \pi  \frac{n}{N}  k \right)}} = \sum\limits_{k = 2}^{6}{1  e^{- j  \left( 2  \pi  \frac{n}{N}  k \right)}}\ \label{eq:10.15}
\end{align}
Die Berechnung muss nicht unbedingt als analytische Rechnung ausgeführt werden, sie kann wegen der endlichen Summe auch numerisch erfolgen. Das Signal und der Betrag des Spektrums sind in Bild \ref{fig:10.1} dargestellt.

\begin{fixedfigure}
  \includegraphics[alt={Rechteckfolge und ihre Diskrete-Fourier-Transformierte}]{new media/12_1_rechteckfolge_und_ihre_diskrete.png}
  \captionof{figure}{Rechteckfolge und ihre Diskrete-Fourier-Transformierte}
  \label{fig:10.1}
\end{fixedfigure}

Die Diskrete-Fourier-Transformierte wird mit dem Spektrum \(X(Ω)\) verglichen, das sich aus der Fourier-Transformation der Folge ergibt.
\begin{align}
X(\Omega) &= \sum\limits_{k = 2}^{6}{1  e^{- j  k  \Omega}} = e^{- j  4  \Omega}  \frac{\sin\left( \frac{5  \Omega}{2} \right)}{\sin\left( \frac{\Omega}{2} \right)} \label{eq:10.16}
\end{align}
Beide Spektren sind in Bild \ref{fig:10.2} dargestellt. Die periodische Fortsetzung der Diskreten-Fourier-Transformierten und die Fourier-Transformierte der Folge sind an den Punkten \(Ω_{n}\) identisch. Allerdings liefert die Diskrete-Fourier-Transformierte nicht den bislang interpretierten Frequenzbereich von - π \ldots{} π sondern das Spektrum im Bereich von 0 \ldots{} 2π. Wegen der Periodizität der Fourier-Transformierten können die Darstellungen jedoch ineinander überführt werden.

\begin{fixedfigure}
  \includegraphics[alt={Spektrum der Rechteckfolge berechnet über Fourier-Transformation\\
und Diskrete-Fourier-Transformation}]{new media/12_2_spektrum_der_rechteckfolge_bere.png}
  \captionof{figure}{Spektrum der Rechteckfolge berechnet über Fourier-Transformation\\
und Diskrete-Fourier-Transformation}
  \label{fig:10.2}
\end{fixedfigure}

\end{example}

\subsection{Grenzen der Diskreten-Fourier-Transformation}\label{sec:grenzen-der-diskreten-fourier-transformation}

Aufgrund der endlichen Beobachtungszeit \(T_{B}\) = NT\textsubscript{A} und des Abtastvorgangs wird das Spektrum mit der Diskreten-Fourier-Transformation nur an endlich vielen Stellen berechnet. Von daher stellt sich die Frage nach der maximalen Frequenz und der Auflösung des berechneten Spektrums.

\subsubsection{Maximale Frequenz des Spektrums der Diskreten-Fourier-Transformation}\label{sec:maximale-frequenz-des-spektrums-der-diskreten-fourier-transformation}

Die maximale Frequenz der Diskreten-Fourier-Transformation ergibt sich aus dem Abtasttheorem. Mit diskreten Abtastwerten eines kontinuierlichen Signals werden nur die Kreisfrequenzen dargestellt, die kleiner sind als die halbe Abtastkreisfrequenz
\begin{align}
\omega \leq \omega_{MAX} &= \frac{\omega_{A}}{2} \label{eq:10.17}
\end{align}
Die maximale Frequenz der Diskreten-Fourier-Transformierten wird deshalb schon mit dem Abtastvorgang festgelegt.

\subsubsection{Spektrale Auflösung der Diskreten-Fourier-Transformation}\label{sec:spektrale-aufluxf6sung-der-diskreten-fourier-transformation}

Der Abstand zweier Frequenzwerte einer Diskreten-Fourier-Transformierten errechnet sich nach Gleichung \eqref{eq:10.14} zu
\begin{align}
\Delta\omega &= \frac{2  \pi}{T_{A}  N} \label{eq:10.18}
\end{align}
Für eine gute spektrale Auflösung, also kleine Frequenzabstände Δω, muss die Anzahl N von Abtastwerten hoch gewählt werden. Wird für die Anzahl N nach Gleichung \eqref{eq:10.18} das Verhältnis von Beobachtungszeit zur Abtastzeit eingesetzt, ergibt sich für die Auflösung
\begin{align}
\Delta\omega &= \frac{2  \pi}{T_{A}  \frac{T_{B}}{T_{A}}} = \frac{2  \pi}{T_{B}} \label{eq:10.19}
\end{align}
Die Auflösung kann demnach nur durch eine Verlängerung der Beobachtungszeit \(T_{B}\) verbessert werden.

\begin{example}{Interpretation der Impulsantwort eines Systems 2. Ordnung mit komplexem Polpaar}
\label{beispiel-interpretation-der-impulsantwort-eines-systems-2.-ordnung-mit-komplexem-polpaar}

An dem Beispiel der Impulsantwort eines Verzögerungsglieds zweiter Ordnung wird der praktische Nutzen der Diskreten-Fourier-Transformation verdeutlicht. Das Signal wird durch Abtastung der Impulsantwort eines PT2-Glieds mit der Übertragungsfunktion
\begin{align}
G(s) &= \frac{1}{1 + 2  d  T  s + T^{2}  s^{2}} \label{eq:10.20}
\end{align}
mit \(T = \SI{31}{\milli\second}\) und \(d = 0.06\) erzeugt. Bild \ref{fig:10.3} zeigt die zeitkontinuierliche Impulsantwort.

\begin{fixedfigure}
  \includegraphics[alt={Darstellung der Impulsantwort eines Systems 2. Ordnung mit komplexem Polpaar}]{new media/12_3_darstellung_der_impulsantwort_e.png}
  \captionof{figure}{Darstellung der Impulsantwort eines Systems 2. Ordnung mit komplexem Polpaar}
  \label{fig:10.3}
\end{fixedfigure}

Die Impulsantwort wird für eine Dauer von \SI{2}{\second} mit \(T_{A} = \SI{1}{\milli\second}\) abgetastet. Die entsprechenden Folgenwerte \(x[k]\) werden mit Hilfe der Diskreten-Fourier-Transformation transformiert. Das Ergebnis ist die Diskrete-Fourier-Transformierte \(X_{n}\). Ihr Betrag ist in Bild \ref{fig:10.4} dargestellt.

\begin{fixedfigure}
  \includegraphics[alt={Darstellung der Diskreten-Fourier-Transformierten für die Impulsantwort eines PT2-Glieds}]{new media/12_4_darstellung_der_diskretenfourie.png}
  \captionof{figure}{Darstellung der Diskreten-Fourier-Transformierten für die Impulsantwort eines PT2-Glieds}
  \label{fig:10.4}
\end{fixedfigure}

Die Signale werden mit einer Abtastzeit von \(T_{A} = \SI{0.001}{\second}\) abgetastet. Wegen der Symmetrie des Spektrums ist die maximal auswertbare Frequenz die halbe Abtastfrequenz
\begin{align}
\frac{\omega_{A}}{2} &= \frac{\pi}{T_{A}} = 3141\ rad/s \label{eq:10.21}
\end{align}
Bild \ref{fig:10.5} zeigt den Amplituden- und Phasengang in dem Spektralbereich von 0 \ldots{} \(ω_{A}\)/2 in logarithmischer Darstellung und zum Vergleich das Spektrum der zeitkontinuierlichen Funktion \(g(t)\). Dabei wird das über die DFT berechnete Spektrum mit der Abtastzeit \(T_{A}\) multipliziert, um die aus der Abtastung resultierende Skalierung zu kompensieren.

\begin{fixedfigure}
  \includegraphics[alt={Spektrum der Impulsantwort eines PT2-Glieds,\\
Vergleich von analytischer Rechnung und Diskreter-Fourier-Transformierter}]{new media/12_5_spektrum_der_impulsantwort_eine.png}
  \captionof{figure}{Spektrum der Impulsantwort eines PT2-Glieds,\\
Vergleich von analytischer Rechnung und Diskreter-Fourier-Transformierter}
  \label{fig:10.5}
\end{fixedfigure}

Die Amplitudengänge stimmen in einem weiten Frequenzbereich überein, nur bei Frequenzen nahe der halben Abtastfrequenz ergeben sich Abweichungen aufgrund der periodischen Wiederholung des Spektrums durch den Abtastvorgang. Die Phasengänge des abgetasteten und des zeitkontinuierlichen Signals sind praktisch identisch.

\end{example}

Bei den beiden bislang diskutierten Beispielen stimmen Fourier-Transformierte der Signalfolge beziehungsweise die Fourier-Transformierte der Impulsantwort und die entsprechende Diskrete-Fourier-Transformierte gut überein. Diese gute Übereinstimmung wird nur erreicht, wenn einige Bedingungen an das Signal und die Wahl des Beobachtungszeitraums erfüllt sind. Auf diese Randbedingungen wird in Abschnitt \ref{sec:10.3} näher eingegangen.

\subsection{Inverse Diskrete-Fourier-Transformierte}\label{sec:inverse-diskrete-fourier-transformierte}

Zu der Diskreten-Fourier-Transformation existiert eine inverse Diskrete-Fourier-Transformation, mit der aus den Werten \(X_{n}\) die abgetastete Folge \(x[k]\) berechnet werden kann.
\begin{align}
x[ k] &= \frac{1}{N}  \sum\limits_{n = 0}^{N - 1}X_{n}  e^{j  \left( 2  \pi  \frac{n}{N}  k \right)} \label{eq:10.22}
\end{align}
Zum Beweis dieser Umkehrformel wird in Gleichung \eqref{eq:10.22} die Gleichung zur Bestimmung der Fourier-Koeffizienten \(X_{n}\)
\begin{align}
X_{n} &= \sum\limits_{k = 0}^{N - 1}{x[ k]  e^{- j  \left( 2  \pi  \frac{n}{N}  k \right)}}\ \label{eq:10.23}
\end{align}
eingesetzt. Daraus ergibt sich
\begin{align}
x[ k] &= \frac{1}{N}  \sum\limits_{n = 0}^{N - 1}X_{n}  e^{j  2  \pi  \frac{n}{N}  k} = \frac{1}{N}  \sum\limits_{n = 0}^{N - 1}{\sum\limits_{m = 0}^{N - 1}{x[ m]}  e^{- j  2  \pi  \frac{n}{N}  m}}  e^{j  2  \pi  \frac{n}{N}  k} = \frac{1}{N}  \sum\limits_{m = 0}^{N - 1}{x[ m]}  \sum\limits_{n = 0}^{N - 1}e^{j  2  \pi  \frac{n}{N}  (k - m)} \label{eq:10.24}
\end{align}
Die Summe über die komplexen Exponentialfunktionen ist nur dann von null verschieden, wenn\\
\(k = m\) ist. Für \(k = m\) wird der Ausdruck
\begin{align}
e^{j  2  \pi  \frac{n}{N}  (k - m)} &= 1 \label{eq:10.25}
\end{align}
N-mal aufsummiert, das Ergebnis beträgt damit N. Für ein bestimmtes k existiert damit nur ein Element, das nicht null wird. Damit ergibt sich
\begin{align}
x[ k] &= \frac{1}{N}  \sum\limits_{m = 0}^{N - 1}{x[ m]}  \sum\limits_{n = 0}^{N - 1}e^{j  2  \pi  \frac{n}{N}  (k - m)} = \frac{1}{N}  x[ k]  N = x[ k] \label{eq:10.26}
\end{align}
Anhand der Gleichung \eqref{eq:10.22} zur Rücktransformation wird deutlich, dass die zu dem diskreten Spektrum \(X_{n}\) gehörende Folge \(x[k]\) ebenfalls periodisch in N ist.
\begin{align}
x[ k + N] &= \frac{1}{N}  \sum\limits_{n = 0}^{N - 1}X_{n}  e^{j  \left( 2  \pi  \frac{n}{N}  (k + N) \right)} = \frac{1}{N}  \sum\limits_{n = 0}^{N - 1}X_{n}  e^{j  \left( 2  \pi  \frac{n}{N}  k \right)}  e^{j  \left( 2  \pi  \frac{n}{N}  N \right)} = \frac{1}{N}  \sum\limits_{n = 0}^{N - 1}X_{n}  e^{j  \left( 2  \pi  \frac{n}{N}  k \right)} = x[ k] \label{eq:10.27}
\end{align}
Die Folge \(x[k]\) wird bei der Rekonstruktion über den Beobachtungszeitraum periodisch fortgesetzt.

\subsection{Fourier-Reihe des periodisch wiederholten Signals}\label{sec:fourier-reihe-des-periodisch-wiederholten-signals}

Bei der Herleitung der Diskreten-Fourier-Transformation in Abschnitt \ref{sec:10.1.1} wird von einem ideal abgetasteten, zeitlich begrenzten Signal ausgegangen. Es wird gezeigt, dass die Diskrete-Fourier-Transformierte \(X_{n}\) bis auf Skalierungsfaktoren den Fourier-Koeffizienten des ideal abgetasteten Signals entspricht.

Die Diskrete-Fourier-Transformation wird aber auch zur numerischen Näherung der Fourier-Reihe des zeitkontinuierlichen Signals verwendet. Die Näherung ergibt sich dabei aus dem Abtastvorgang. Bei der idealen Abtastung des Signals wird das Spektrum des zeitkontinuierlichen Signals periodisch in \(ω_{A}\) wiederholt und mit 1/\(T_{A}\) skaliert. Das gefensterte Signal ist zeitbegrenzt und damit nicht bandbegrenzt. Deshalb vermischt sich das Spektrum im Basisband mit seinen periodischen Wiederholungen. Die Diskrete-Fourier-Transformation ist damit nur eine Näherung der Fourier-Reihe. Je schneller das Spektrum abfällt und je höher die Abtastfrequenz \(ω_{A}\) ist, desto besser ist die Näherung.
\begin{align}
A_{n} \approx \frac{1}{N_{A}}  X_{n} \label{eq:10.28}
\end{align}
Der Faktor 1/N ergibt sich aus der Herleitung in Abschnitt \ref{sec:10.1.1}, der Faktor 1/\(T_{A}\) ergibt sich aus dem Abtastvorgang.

\subsection{Signalformen und ihr Spektrum}\label{sec:signalformen-und-ihr-spektrum}

Der Zusammenhang zwischen kontinuierlichen und diskreten Signalen und ihren Spektren wird anhand eines dreieckförmigen Signals \(x(t)\) verdeutlicht.
\begin{align}
x(t) &= (t + 1)  \sigma(t + 1) - 2  t  \sigma(t) + (t - 1)  \sigma(t - 1) \label{eq:10.29}
\end{align}
Die Fourier-Transformierte des Signals \(x(t)\) ergibt sich mit den bekannten Korrespondenzen der Fourier-Transformation und \(T = 1\) zu
\begin{align}
X(\omega) &= \left( \frac{\sin\left( \frac{\omega}{2} \right)}{\frac{\omega}{2}} \right)^{2} = \sin^{2}\left( \frac{\omega}{2} \right)  \frac{4}{\omega^{2}} \label{eq:10.30}
\end{align}
Das zeitkontinuierliche Signal \(x(t)\) und das Spektrum \(X(ω)\) sind in Bild \ref{fig:10.6} dargestellt.

\begin{figure}
  \centering
  \includegraphics[alt={Dreieck-Funktion und das Spektrum für das zeitkontinuierliche Signal}]{new media/12_6_dreieckfunktion_und_das_spektru.png}
  \caption{Dreieck-Funktion und das Spektrum für das zeitkontinuierliche Signal}
  \label{fig:10.6}
\end{figure}

Das Signal \(x(t)\) wird in Abständen von \(T_{A} = 0.2\) ideal abgetastet. Es entsteht das Signal \(x_{A}(t)\).
\begin{align}
x_{A}(t) &= \sum\limits_{k = - \infty}^{\infty}{x\left( k  T_{A} \right)  \delta\left( t - k  T_{A} \right)} \label{eq:10.31}
\end{align}
Nach idealer Abtastung des Signals im Zeitbereich liegen die Funktionswerte nur noch an den Stellen k·\(T_{A}\) vor. Durch die Diskretisierung wird das Signal im Frequenzbereich in \(ω_{A}\) = 2π/\(T_{A}\) = 10π periodisch fortgesetzt und mit einem Faktor 1/\(T_{A} = 5\) multipliziert. Damit ergibt sich das Spektrum \(X_{A}(ω)\) des abgetasteten Signals.
\begin{align}
X_{A}(\omega) &= \frac{1}{T_{A}}  \sum\limits_{n = - \infty}^{\infty}{X\left( \omega - n  \omega_{A} \right)} = 5  \sum\limits_{n = - \infty}^{\infty}\frac{4}{(\omega - n  10  \pi)^{2}}  \sin^{2}\left( \frac{\omega - n  10  \pi}{2} \right) \label{eq:10.32}
\end{align}
Das ideal abgetastete Signal und das Spektrum \(X_{A}(ω)\) sind in Bild \ref{fig:10.7} dargestellt.

\begin{figure}
  \centering
  \includegraphics[alt={Ideal abgetastete Dreieckfunktion und ihr Spektrum, Spektrum wird mit 1/\(T_{A}\) und periodisch fortgesetzt}]{new media/12_7_ideal_abgetastete_dreieckfunkti.png}
  \caption{Ideal abgetastete Dreieckfunktion und ihr Spektrum, Spektrum wird mit 1/\(T_{A}\) und periodisch fortgesetzt}
  \label{fig:10.7}
\end{figure}

Die Folge \(x_{D}\){[}k{]} ergibt sich aus den Gewichten des ideal abgetasteten Signals x(kT\textsubscript{A}).
\begin{align}
x_{D}[ k] &= x\left( k  T_{A} \right) \label{eq:10.33}
\end{align}
Die Fourier-Transformierte \(X_{D}(Ω)\) entspricht grundsätzlich dem Spektrum \(X_{A}(ω)\), allerdings wird es als Funktion der normierten Kreisfrequenz Ω = ωT\textsubscript{A} dargestellt. Mit der Substitution \(ω = Ω\)/\(T_{A}\) ergibt sich
\begin{align}
X_{D}(\Omega) &= \left. \ \frac{1}{T_{A}}  \sum\limits_{n = - \infty}^{\infty}{X\left( \omega - n  \omega_{A} \right)} \right|_{\omega = \frac{\Omega}{T_{A}}} = 5  \sum\limits_{n = - \infty}^{\infty}\frac{4}{(5  \Omega - n  10  \pi)^{2}}  \sin^{2}\left( \frac{5  \Omega - n  10  \pi}{2} \right) \label{eq:10.34}
\end{align}
Die diskrete Signalfolge und das Spektrum \(X_{D}(Ω)\) sind in Bild \ref{fig:10.8} dargestellt.

\begin{figure}
  \centering
  \includegraphics[alt={Dreieckfolge und ihr Spektrum, Spektrum wird als Funktion der normierten Kreisfrequenz dargestellt}]{new media/12_8_dreieckfolge_und_ihr_spektrum_s.png}
  \caption{Dreieckfolge und ihr Spektrum, Spektrum wird als Funktion der normierten Kreisfrequenz dargestellt}
  \label{fig:10.8}
\end{figure}

Die Signalfolge \(x_{D}\){[}k{]} wird über einen Zeitraum - 5 ≤ k ≤ 4 beobachtet. Mit den entsprechenden Werten \(x_{W}\){[}k{]} der Signalfolge wird eine Diskrete-Fourier-Transformierte \(X_{n}\) berechnet. Es werden \(N = 10\) Abtastwerte verwendet. Deshalb berechnen sich die Frequenzen \(Ω_{n}\) mit \(n = 0\) ... 9 zu
\begin{align}
\Omega_{n} &= \frac{2  \pi}{N}  n = \frac{2  \pi}{10}  n \label{eq:10.35}
\end{align}
Da das Signal \(x[k]\) sowieso zeitbegrenzt war, wirkt sich die Fensterung nicht auf das Spektrum aus. Das Spektrum der Signalfolge \(x_{D}\){[}k{]} entspricht damit an den Stellen \(Ω = Ω_{n}\) den Werten der diskreten Fourier-Transformation X(\(Ω_{n}\)) = \(X_{n}\). Wegen der Abtastung im Frequenzbereich wird das Signal im Zeitbereich periodisch in N wiederholt. Die periodisch fortgesetzte Signalfolge und die Diskrete-Fourier-Transformierte sind in Bild \ref{fig:10.9} dargestellt.

\begin{figure}
  \centering
  \includegraphics[alt={Dreieckfolge und ihre DFT (\(N = 10\)), Zeitsignal periodisch in N fortgesetzt}]{new media/12_9_dreieckfolge_und_ihre_dft_n_10.png}
  \caption{Dreieckfolge und ihre DFT (\(N = 10\)), Zeitsignal periodisch in N fortgesetzt}
  \label{fig:10.9}
\end{figure}

Wird der Ausschnitt des zeitkontinuierlichen Signals für - 1 ≤ t \textless{} 1 periodisch mit \(T_{P} = 2\) wiederholt, kann das Spektrum als komplexen Fourier-Reihe berechnet werden. Die Fourier-Koeffizienten errechnet sich zu
\begin{align}
A_{n} &= \frac{1}{T_{P}}  \int_{- T_{P}/2}^{T_{P}/2}{x(t)  e^{- j  n  \omega_{0}  t}dt} = \frac{1}{2}  \left( \int_{- 1}^{0}{(1 + t)  e^{- j  n  \omega_{0}  t}dt} + \int_{0}^{1}{(1 - t)  e^{- j  n  \omega_{0}  t}dt} \right) \label{eq:10.36}
\end{align}
Die Gleichung stimmt bis auf den Faktor 1/2 und ω = nω\textsubscript{0} mit der Bestimmungsgleichung der Fourier-Transformation überein. Damit können die Fourier-Koeffizienten angegeben werden als
\begin{align}
A_{n} &= \frac{1}{T_{P}}  \left. \ X(\omega) \right|_{\omega = n  \omega_{0}} = \frac{1}{2}  \sin^{2}\left( \frac{n  \omega_{0}}{2} \right)  \frac{4}{n^{2}  {\omega_{0}}^{2}} \label{eq:10.37}
\end{align}
Sie sind in Bild \ref{fig:10.10} dargestellt.

\begin{figure}
  \centering
  \includegraphics[alt={Periodisch fortgesetzte Dreieck-Funktion und ihre Fourier-Koeffizienten}]{new media/12_10_periodisch_fortgesetzte_dreiec.png}
  \caption{Periodisch fortgesetzte Dreieck-Funktion und ihre Fourier-Koeffizienten}
  \label{fig:10.10}
\end{figure}

Mit den Ausführungen in Abschnitt gilt zwischen Fourier-Reihe und diskreter Fourier-Transformation näherungsweise
\begin{align}
A_{n} \approx \frac{1}{N}  X_{n} &= \frac{1}{10}  X_{n} \label{eq:10.38}
\end{align}
Da das Spektrum des zeitkontinuierlichen Signals \(x_{p}(t)\) nicht bandbegrenzt und die Abtastfrequenz gering ist, ist diese Näherung nur im niedrigen Frequenzbereich ω \textless\textless{} \(ω_{A}\)/2 sinnvoll. Tabelle \ref{tab:10.1} vergleicht die DFT der Abtastwerte und Fourier-Reihe des periodisch wiederholten Signalausschnitts numerisch.

\begin{table}
  \centering
  \caption{Vergleich DFT der Abtastwerte und Fourier-Reihe des periodisch wiederholten Signalausschnitts \(x_{P}(t)\)}
  \label{tab:10.1}
  \begin{tabular}{llllllllllll}
    \toprule
    Index n & - 5 & - 4 & - 3 & - 2 & - 1 & 0 & 1 & 2 & 3 & 4 & 5 \\
    \midrule
    Frequenz \(ω_{n}\) / rad/s & 15.71 & 12.56 & 9.424 & 6.283 & 3.141 & 0 & 3.141 & 6.283 & 9.424 & 12.56 & 15.71 \\
    DFT von \(x_{W}\){[}k{]} & 0.192 & 0.000 & 0.298 & 0.000 & 2.088 & 5.000 & 2.088 & 0.000 & 0.298 & 0.000 & 0.192 \\
    NAn von \(x_{p}(t)\) & 0.081 & 0.000 & 0.225 & 0.000 & 2.028 & 5.000 & 2.028 & 0.000 & 0.225 & 0.000 & 0.081 \\
    \bottomrule
  \end{tabular}
\end{table}

Das Beispiel zeigt, dass aus der Diskretisierung im Zeitbereich bei der Abtastung von Signalen eine periodische Wiederholung der Spektren im Frequenzbereich folgt. Umgekehrt führt eine Diskretisierung des Spektrums zu einer periodischen Wiederholung der abgetasteten Folge. Die Diskrete-Fourier-Transformation berechnet demnach nicht das Spektrum der Signalfolge im definierten Zeitfenster, sondern das Spektrum der periodisch fortgesetzten Signalfolge. Die sich daraus ergebenden Konsequenzen für nicht periodische Signale werden in Abschnitt \ref{sec:10.3} aufgegriffen.

\subsection{Zusammenfassung Diskrete-Fourier-Transformation}\label{sec:zusammenfassung-diskrete-fourier-transformation}

Abschließend fasst Tabelle \ref{tab:10.2} die Eigenschaften der Diskreten-Fourier-Transformation zusammen.

\begin{table}
  \centering
  \caption{Übersicht über die Diskrete-Fourier-Transformation}
  \label{tab:10.2}
  \begin{tabular}{ll}
    \toprule
    Eigenschaft & Mathematische Beschreibung \\
    \midrule
    Diskrete-Fourier-Transformation & \(X_{n} = \sum\limits_{k = 0}^{N - 1}{x\lbrack k\rbrack  e^{- j  \left( 2  \pi  \frac{n}{N}  k \right)}}\ \) \\
    Diskrete-Fourier-Transformation und 
Fourier-Transformation von Signalfolgen
 & \(X_{n} = X\left( \Omega_{n} \right)\) \\
    Frequenzen & \(\omega_{n} = \frac{\Omega_{n}}{T_{A}} = \frac{2  \pi  n}{T_{A}  N}\) \\
    Maximale Frequenz & \(\omega_{MAX} = \frac{\omega_{A}}{2}\) \\
    Spektrale Auflösung & \(\Delta\omega = \frac{2  \pi}{T_{A}  N} = \frac{2  \pi}{T_{B}}\) \\
    Inverse Diskrete-Fourier-Transformation & \(x\lbrack k\rbrack = \frac{1}{N}  \sum\limits_{n = 0}^{N - 1}X_{n}  e^{j  \left( 2  \pi  \frac{n}{N}  k \right)}\) \\
    Fourier-Reihe des ideal abgetasteten und 
periodisch wiederholten Signals \(x_{A}(t)\)
 & \(A_{n} = \frac{1}{N  T_{A}}  X_{n}\) \\
    Fourier-Reihe des periodisch wiederholten 
Signalausschnitts \(x_{P}(t)\)
 & \(A_{n} \approx \frac{1}{N}  X_{n}\) \\
    \bottomrule
  \end{tabular}
\end{table}

\section{Eigenschaften der Diskreten-Fourier-Transformation}\label{sec:10.2}

Die Diskrete-Fourier-Transformation hat Eigenschaften, die zum großen Teil vergleichbar zur Fourier-Transformation für Folgen sind. Aufgrund der periodischen Fortsetzung der Signalfolge über das Beobachtungsfenster hinaus ergeben sich teilweise neue Aspekte. Da die Diskrete-Fourier-Transformation typischerweise nicht analytisch bestimmt wird, sind die Rechenregeln nur zum allgemeinen Verständnis und für die Herleitung effizienter Algorithmen wie der Fast-Fourier-Transformation von Bedeutung.

\subsection{Linearität}\label{sec:linearituxe4t-4}

Die Diskrete-Fourier-Transformierte einer Linearkombination von zwei Signalfolgen \(x_{1}\){[}k{]} und \(x_{2}\){[}k{]} gleicher Länge ergibt sich aus derselben Linearkombination \(X_{1n}\) und \(X_{2n}\) der jeweiligen Diskreten-Fourier-Transformierten \(X_{1n}\) und \(X_{2n}\).
\begin{align}
{X_{n} } &=  {\sum\limits _{k=0}^{N-1}\left(\nu _{1}  x_{1} \left[k\right]+\nu _{2}  x_{2} \left[k\right]\right) e^{-j \left(2 \pi  \frac{n}{N}  k\right)}  =\sum\limits _{k=0}^{N-1}\nu _{1}  x_{1} \left[k\right] e^{-j \left(2 \pi  \frac{n}{N}  k\right)}  +\sum\limits _{k=0}^{N-1}\nu _{2}  x_{2} \left[k\right] e^{-j \left(2 \pi  \frac{n}{N}  k\right)}  } \\
{} &=  {\nu _{1}  X_{1n} +\nu _{2}  X_{2n} }
 \label{eq:10.39}
\end{align}
Wie bei der Fourier-Transformation von Signalfolgen gilt also das Linearitätsprinzip.

\subsection{Symmetrie der Diskreten-Fourier-Transformierten reeller Signalfolgen}\label{sec:10.2.2}

Die Diskrete-Fourier-Transformierte reeller Folgen hat eine konjugiert komplexe Symmetrie. Die Symmetrie kann durch Umformen der Definitionsgleichung aufgezeigt werden.
\begin{align}
X_{N - n} &= \sum\limits_{k = 0}^{N - 1}{x[ k]  e^{- j  \left( 2  \pi  \frac{N - n}{N}  k \right)}} = \sum\limits_{k = 0}^{N - 1}{x[ k]  e^{- j  \left( 2  \pi  \frac{N}{N}  k \right)}  e^{- j  \left( 2  \pi  \frac{- n}{N}  k \right)}} = \sum\limits_{k = 0}^{N - 1}{x[ k]  e^{j  \left( 2  \pi  \frac{n}{N}  k \right)}}\text{=}\text{X}_{n}^{*}\ \label{eq:10.40}
\end{align}
Die übrigen Symmetrieregeln der Fourier-Transformation von Signalfolgen gelten sinngemäß. Sie sind in Tabelle \ref{tab:10.3} zusammengefasst.

\begin{table}
  \centering
  \caption{Symmetrieregeln der Fourier-Transformation für Signalfolgen}
  \label{tab:10.3}
  \begin{tabular}{ll}
    \toprule
    Folge \(x[k]\) & Diskrete-Fourier-Transformierte \(X_{n}\) \\
    \midrule
    Gerade Folge \(x[k]\) & Reelles Spektrum \(X_{n}\) \\
    Ungerade Folge \(x[k]\) & Imaginäres Spektrum \(X_{n}\) \\
    x*{[}k{]} & X*\textsubscript{-n} \\
    x*{[}-k{]} & X*\textsubscript{n} \\
    \bottomrule
  \end{tabular}
\end{table}

\subsection{Dualität der Diskreten-Fourier-Transformation}\label{sec:dualituxe4t-der-diskreten-fourier-transformation}

Die Definitionsgleichung der Diskreten-Fourier-Transformation
\begin{align}
X_{n} &= \sum\limits_{k = 0}^{N - 1}{x[ k]  e^{- j  \left( 2  \pi  \frac{n}{N}  k \right)}}\ \label{eq:10.41}
\end{align}
und der inversen Diskreten-Fourier-Transformation
\begin{align}
x[ k] &= \frac{1}{N}  \sum\limits_{n = 0}^{N - 1}X_{n}  e^{j  \left( 2  \pi  \frac{n}{N}  k \right)} \label{eq:10.42}
\end{align}
sind sich mathematisch sehr ähnlich. Sie unterscheiden sich um einen Faktor 1/N und im Vorzeichen des Exponenten der Exponentialfunktion. Aus dieser Symmetrie ergeben sich weitere Interpretationsmöglichkeiten. Zum Beispiel errechnet sich bei bekannter periodischer Diskrete-Fourier-Transformierte \(X_{n}\) einer periodischen Signalfolge \(x[k]\) das Spektrum der Signalfolge \(y[k]\) = \(X_{k}\) zu
\begin{align}
Y_{n} &= \sum\limits_{k = 0}^{N - 1}{y[ k]  e^{- j  \left( 2  \pi  \frac{n}{N}  k \right)}} = N  x[ - k] \label{eq:10.43}
\end{align}
Statt eines Beweises wird die Dualität an einem Beispiel illustriert.

\begin{example}{Dualität der Diskreten-Fourier-Transformation}
\label{beispiel-dualituxe4t-der-diskreten-fourier-transformation}

Eine in N periodisch wiederholten Impulsfolge \(δ[k]\) kann als Summe dargestellt werden.
\begin{align}
x[ k] &= \sum\limits_{n = - \infty}^{\infty}{\delta[ k - n  N]}\ \label{eq:10.44}
\end{align}
Die Diskrete-Fourier-Transformierte wird über einen Signalausschnitt der Länge N berechnet.
\begin{align}
X_{n} &= \sum\limits_{k = 0}^{N - 1}{x[ k]  e^{- j  \left( 2  \pi  \frac{n}{N}  k \right)}}\  = \sum\limits_{k = 0}^{N - 1}{\delta[ k]  e^{- j  \left( 2  \pi  \frac{n}{N}  k \right)}}\  = 1\ \label{eq:10.45}
\end{align}
Mit dieser Korrespondenz und der Dualität kann die Diskrete-Fourier-Transformierte der Folge\\
\(y[k]\) = 1 berechnet werden. Es gilt der Zusammenhang
\begin{align}
y[ k] &= X_{k} = 1\ \label{eq:10.46}
\end{align}
Mit der Dualität der Diskreten-Fourier-Transformierten ergibt sich
\begin{align}
Y_{n} &= N  x[ - k] = N  \sum\limits_{n = - \infty}^{\infty}{\delta\left[ - (k - n  N) \right]}\  = N  \sum\limits_{n = - \infty}^{\infty}{\delta[ - k + n  N]} \label{eq:10.47}
\end{align}
Bild \ref{fig:10.11} stellt die Diskreten-Fourier-Transformierten für die periodische Impuls- und die Einsfolge als Beispiel für Dualität dar. Es werden \(N = 16\) Abtastwerte ausgewertet.

\begin{fixedfigure}
  \includegraphics[alt={Darstellung der Diskreten-Fourier-Transformierten für Impuls- und Einsfolge}]{new media/12_11_darstellung_der_diskretenfouri.png}
  \captionof{figure}{Darstellung der Diskreten-Fourier-Transformierten für Impuls- und Einsfolge}
  \label{fig:10.11}
\end{fixedfigure}

\end{example}

Diese Dualität existiert in analoger Weise für die Fourier-Transformation zeitkontinuierlicher Signale. Im Gegensatz zur Diskreten-Fourier-Transformation weist die Fourier-Transformation von Signalfolgen dieses Dualitätsprinzip nicht auf, da das Signal \(x[k]\) eine zeitdiskrete Folge und das Spektrum \(X(Ω)\) eine kontinuierliche Funktion ist.

\subsection{Zeitverschiebung}\label{sec:zeitverschiebung}

Die Diskrete-Fourier-Transformierte einer verschobenen Folge x{[}k - \(k_{0}\){]} kann aus der Diskreten-Fourier-Transformation der nicht verschobenen Folge berechnet werden.
\begin{align}
{\sum\limits _{k=0}^{N-1}x\left[k-k_{0} \right] e^{-j \left(2 \pi  \frac{n}{N}  k\right)}  } &=  {\sum\limits _{k=0}^{N-1}x\left[k-k_{0} \right] e^{-j \left(2 \pi  \frac{n}{N}  \left(k-k_{0} \right)\right)}  e^{-j \left(2 \pi  \frac{n}{N}  k_{0} \right)}  } \\
{} &=  {e^{-j \left(2 \pi  \frac{n}{N}  k_{0} \right)}  \sum\limits _{k=0}^{N-1}x\left[k\right] e^{-j \left(2 \pi  \frac{n}{N}  k\right)}  =e^{-j \left(2 \pi  \frac{n}{N}  k_{0} \right)}  X_{n} {\; }}
 \label{eq:10.48}
\end{align}
Sinngemäß entsprechen diese Regeln der Verschiebungsregel der Fourier-Transformation für Signalfolgen. Zu beachten ist, dass \(x[k]\) nicht die Folge in der Beobachtungszeit \(T_{B}\) ist, sondern die periodische Fortsetzung dieser Folge. Wegen der periodischen Fortsetzung wird die Zeitverschiebung auch als zirkuläre Verschiebung der Folge bezeichnet.

\subsection{Faltung und periodische Faltung}\label{sec:faltung-und-periodische-faltung}

Bei den bislang diskutierten Transformationen entspricht die Faltung zweier Signale im Zeitbereich der Multiplikation der Fourier-Transformierten im Frequenzbereich. Diese Eigenschaft gilt auch für die Diskrete-Fourier-Transformation. Allerdings muss auch hier berücksichtigt werden, dass die Signale bei der Diskreten-Fourier-Transformation periodisch fortgesetzt werden. Die Auswirkung der periodischen Fortsetzung und Möglichkeiten zur Abhilfe werden an einem Beispiel diskutiert.

\begin{example}{Periodische Faltung der Diskreten-Fourier-Transformation}
\label{beispiel-periodische-faltung-der-diskreten-fourier-transformation}

Gegeben sind die beiden untenstehenden Folgen \(x_{1}\){[}k{]} und \(x_{2}\){[}k{]}.

\begin{fixedfigure}
  \includegraphics[alt={Signalfolgen als Beispiel für eine periodische Faltung}]{new media/12_12_signalfolgen_als_beispiel_fr_e.png}
  \captionof{figure}{Signalfolgen als Beispiel für eine periodische Faltung}
\end{fixedfigure}

Die Faltung der beiden Signale
\begin{align}
y[ k] &= x_{1}[ k]*x_{2}[ k] \label{eq:10.49}
\end{align}
kann zum Beispiel als grafische Faltung ausgeführt werden. Für das Signal \(y[k]\) ergeben sich die in Tabelle \ref{tab:10.4} dargestellten Werte.

\begin{center}
  \captionof{table}{Ergebnis der Faltung der Signalfolgen \(x_{1}\){[}k{]} und \(x_{2}\){[}k{]} im Zeitbereich}
  \label{tab:10.4}
  \begin{tabular}{llllllll}
    \toprule
    k & 0 & 1 & 2 & 3 & 4 & 5 & 6 \\
    \midrule
    \(y[k]\) & 4 & 14 & 17 & 41 & 40 & 31 & 40 \\
    \(y_{DFT}\){[}k{]} & 44 & 45 & 57 & 41 & & & \\
    \bottomrule
  \end{tabular}
\end{center}

Zur Berechnung der Faltung mit der Diskreten-Fourier-Transformation müssen die beiden Diskreten-Fourier-Transformierten berechnet werden. Die Diskrete-Fourier-Transformierte des Ausgangssignals ergibt sich aus dem Produkt von \(X_{1n}\) und \(X_{2n}\). Tabelle \ref{tab:10.5} stellt die Werte dar.

\begin{center}
  \captionof{table}{Diskrete-Fourier-Transformierte \(X_{1}\) und \(X_{2}\) der Signalfolgen \(x_{1}\){[}k{]} und \(x_{2}\){[}k{]}}
  \label{tab:10.5}
  \begin{tabular}{lllll}
    \toprule
    n & 0 & 1 & 2 & 3 \\
    \midrule
    \(X_{1n}\) & 11 & - 1 + 2j & - 5 & - 1 - 2j \\
    \(X_{2n}\) & 17 & 1 + 6j & - 3 & 1 - 6j \\
    \(Y_{n}\) & 187 & - 13 - 4j & 15 & - 13 + 4j \\
    \bottomrule
  \end{tabular}
\end{center}

Durch Rücktransformation ergibt sich das Signal \(y_{DFT}\){[}k{]}, das bereits in Tabelle \ref{tab:10.4} eingetragen ist. Wie deutlich zu erkennen ist, stimmt das Ergebnis der beiden Verfahren nicht überein. Das liegt daran, dass bei der Berechnung der DFT implizit angenommen wird, dass das Signal im Zeitbereich periodisch fortgesetzt wird. Die über die Diskrete-Fourier-Transformation berechnete Faltung wird deshalb auch als periodische Faltung bezeichnet.

Bei der Faltung im Zeitbereich wird das Signal außerhalb der Messpunkte zu null angenommen. Dieser Rechenvorgang nennt sich aperiodische Faltung. Für periodisch fortgesetzte Signale wird auch das Ergebnis der Faltung periodisch. Die Periodizität tritt bereits nach vier Werten auf. Daher ist es auch nicht notwendig, mehr als vier Werte für die Faltung über die DFT zu berechnen.

Damit das Ergebnis der periodischen Faltung über die DFT dem Ergebnis der aperiodischen Faltung entspricht, müssen die beiden Signale jeweils mit mindestens N - 1 Nullen aufgefüllt werden, in diesem Fall also mit drei Nullen.

\end{example}

Die unterschiedlichen Arten der Faltungsoperation und ihre Implementierung im Frequenzbereich sind bei dem Einsatz der Diskreten-Fourier-Transformation Signalverarbeitung von wesentlicher Bedeutung und werden in Kapitel \ref{sec:10.6} wieder aufgegriffen.

\section{Spektralschätzung über die Diskreten-Fourier-Transformation}\label{sec:10.3}

Mit der Diskreten-Fourier-Transformation wird einer endlichen Signalfolge \(x_{w}\){[}k{]} ein Spektrum zugewiesen. Die Beispiele in Abschnitt \ref{sec:10.1} zeigen, dass unter günstigen Bedingungen das Spektrum der Signalfolge \(x_{w}\){[}k{]} dem Spektrum des zeitkontinuierlichen Signals \(x(t)\) weitgehend entspricht. Die Güte der Übereinstimmung ist von der Signalverarbeitung abhängig, die zwischen dem zeitkontinuierlichen Signal \(x(t)\) und den Eingangswerten der Diskreten-Fourier-Transformation \(x_{w}\){[}k{]} liegt. Dieser Zusammenhang wird in diesem Abschnitt analysiert.

\subsection{Signalfluss und Einflüsse auf die Spektralschätzung}\label{sec:signalfluss-und-einfluxfcsse-auf-die-spektralschuxe4tzung}

Bild \ref{fig:10.13} zeigt den Signalfluss, der bei der Spektralschätzung verwendet wird, mit den entsprechenden Signalbezeichnungen im Zeit- und Frequenzbereich.

\begin{figure}
  \centering
  \includegraphics[alt={Signalfluss bei der Spektralschätzung}]{new media/12_13_signalfluss_bei_der_spektralsc.pdf}
  \caption{Signalfluss bei der Spektralschätzung}
  \label{fig:10.13}
\end{figure}

Der Signalfluss beginnt mit dem zeitkontinuierlichen Signal \(x(t)\) und dem zu schätzenden Spektrum \(X(ω)\). Damit bei der Abtastung des Signals Aliasing vermieden wird, muss das Signal bandbegrenzt sein. Ist das nicht der Fall, wird das Signal mit einem Anti-Aliasing-Filter bandbegrenzt, damit das Abtasttheorem eingehalten wird. Das bandbegrenzte Signal \(x_{TP}(t)\) weist das Spektrum \(X_{TP}(ω)\) auf.
\begin{align}
X_{TP}(\omega) &= G_{TP}(\omega)  X(\omega) \label{eq:10.50}
\end{align}
Spektralanteile des Signals \(x(t)\), die bei der Tiefpass-Filterung eliminiert werden, stehen nicht mehr zur Verfügung. Der folgende Prozess kann deshalb bestenfalls das Spektrum \(X_{TP}(ω)\) bestimmen. Nur wenn das Signal \(x(t)\) bereits bandbegrenzt ist, sind die Spektren \(X(ω)\) und \(X_{TP}(ω)\) identisch.

Im nächsten Schritt wird das Signal abgetastet, wobei hier von der idealen Abtastung ausgegangen wird. Nach den Ausführungen in Kapitel \hyperref[signalabtastung-und-rekonstruktion]{2} wird durch die Abtastung des Signals das Spektrum periodisch in \(ω_{A}\) wiederholt und mit dem Faktor 1/\(T_{A}\) multipliziert. Nach dem Abtasten ergibt sich damit das Spektrum
\begin{align}
X_{A}(\omega) &= \frac{1}{T_{A}}  \sum\limits_{k = - \infty}^{\infty}{X_{TP}\left( \omega - k  \omega_{A} \right)} \label{eq:10.51}
\end{align}
Im Fall einer nicht idealen Abtastung kann die dadurch entstehende Veränderung des Spektrums durch einen inversen Filter kompensiert werden.

Der Übergang des ideal abgetasteten Signals \(x_{A}(t)\) zur diskreten Signalfolge \(x_{D}\){[}k{]} ist formal. Unter der Bedingung Ω = ωT\textsubscript{A} ist das Spektrum \(X_{A}(ω)\) identisch zu dem Spektrum der Signalfolge \(x_{D}\){[}k{]}, das sich mit der Fourier-Transformation für Signalfolgen berechnet zu
\begin{align}
X_{D}(\Omega) &= \sum\limits_{k = - \infty}^{\infty}{x_{D}[ k]  e^{- j  k  \Omega}} \label{eq:10.52}
\end{align}
Die Signalfolge \(x_{D}\){[}k{]} erstreckt sich im Allgemeinen von - ∞ ≤ k ≤ ∞, sie ist also unendlich lang. Für die Verarbeitung muss sie auf eine endliche Anzahl von N Abtastwerten reduziert werden. Dazu wird sie mit einer Fensterfunktion der Länge \(T_{B}\) multipliziert. Sie ist ein Vielfaches der Abtastzeit \(T_{A}\).
\begin{align}
T_{B} &= N  T_{A} \label{eq:10.53}
\end{align}
Damit liegen von dem Signal N Abtastwerte mit den Indizes \(k = 0\) \ldots{} N - 1 vor, und die Signalfolge ist zeitbegrenzt. Wird als Fensterfunktion ein Rechteckfenster verwendet, ergibt sich
\begin{align}
x_{w}[ k] &= x_{D}[ k]  w[ k] = x_{D}[ k]  \left( \sigma[ k] - \sigma[ k - N] \right) \label{eq:10.54}
\end{align}
Der Multiplikation im Zeitbereich entspricht die Faltung im Frequenzbereich. Durch die Fensterung wird das Spektrum \(X_{D}(Ω)\) mit dem Spektrum der verwendeten Fensterfunktion \(W(Ω)\) gefaltet. Das Spektrum der mit einem Rechteckfenster multiplizierten Signalfolge ergibt sich zu
\begin{align}
X_{w}(\Omega) &= X_{D}(\Omega)*W(\Omega) = X_{D}(\Omega)*\frac{\sin\left( \frac{N  \Omega}{2} \right)}{\sin\left( \frac{\Omega}{2} \right)}  e^{- j  \frac{(N - 1)  \Omega}{2}} \label{eq:10.55}
\end{align}
Die Zeitbegrenzung führt zur Faltung des Spektrums mit dem sogenannten Dirichlet-Kern, was zu einer Veränderung des Spektrums führt. Die Veränderung des Spektrums durch die Faltung mit dem Spektrum der Fensterfunktion wird als Leck-Effekt (Leakage) bezeichnet. Dieser Effekt wird in Abschnitt \ref{sec:10.3.2} wieder aufgegriffen. Da nur N Abtastwerte vorliegen, kann die Fourier-Transformierte \(X_{W}(Ω)\) nur an N Stellen bestimmt werden.
\begin{align}
\Omega_{n} &= \frac{2  \pi  n}{N} \label{eq:10.56}
\end{align}
Das Spektrum \(X_{n}\) stimmt an diesen Stellen mit der Fourier-Transformierten \(X_{W}(Ω)\) der gefensterten Signalfolge \(x_{W}\){[}k{]} überein.
\begin{align}
X_{w}\left( \Omega_{n} \right) &= X_{n} = \sum\limits_{k = 0}^{N - 1}{x_{w}[ k]  e^{- j  \left( 2  \pi  \frac{n}{N}  k \right)}}\ \label{eq:10.57}
\end{align}
Die Herleitung zeigt, dass das berechnete Spektrum \(X_{n}\) aus folgenden Gründen nicht dem Spektrum \(X(ω)\) des Signals \(x(t)\) entspricht:

\begin{itemize}
\item
  Es liegen nur einzelne Werte \(X_{n}\) des Spektrums vor, \(X_{W}(ω)\) ist eine kontinuierliche Funktion. Die kontinuierliche Funktion muss aus den Abtastwerten rekonstruiert werden.
\item
  Das Signal \(x_{w}\){[}k{]} ist gegenüber dem Signal \(x[k]\) zeitbegrenzt. Deshalb wird das ursprüngliche Spektrum mit dem Spektrum der verwendeten Fensterfunktion gefaltet.
\item
  Das Spektrum ist mit dem Faktor \(f_{A} = 1\)/\(T_{A}\) multipliziert. Diese Multiplikation muss rückgängig gemacht werden.
\item
  Das Signal \(x(t)\) muss für die Abtastung ggf. bandbegrenzt werden.
\end{itemize}

Das berechnete Spektrum T\textsubscript{A}\(X_{n}\) ist demnach eine Schätzung oder Approximation des zu bestimmenden Spektrums \(X(ω)\). Daraus ergibt sich die Bezeichnung Spektralschätzung.

\subsection{Auswirkung einer Fensterung von Signalen -- Leakage}\label{sec:10.3.2}

Bei der Fourier-Transformation wird eine Signalfolge unendlich lange beobachtet. Im Gegensatz dazu wird bei der Diskreten-Fourier-Transformation die Signalfolge für einen endlichen Zeitraum beobachtet. Mathematisch kann die Beobachtung des Signals für einen endlichen Zeitraum über die Multiplikation mit einer rechteckförmigen Fensterfunktion \(w[k]\) beschrieben werden.
\begin{align}
w[ k] &= \sigma[ k] - \sigma[ k - N] \label{eq:10.58}
\end{align}
Das zeitlich begrenzte Signal hat damit die Beschreibung
\begin{align}
x_{W}[ k] &= \ x_{D}[ k]  w[ k] \label{eq:10.59}
\end{align}
Nach den Rechenregeln zur Fourier-Transformation berechnet sich das Spektrum \(X_{W}(Ω)\) des zeitlich begrenzten Signals \(x_{W}\){[}k{]} aus der Faltung der Spektren \(X_{D}(Ω)\) und \(W(Ω)\). Das Spektrum \(W(Ω)\) wird als Dirichlet-Kern bezeichnet und berechnet sich zu
\begin{align}
W(\Omega) &= e^{- j  \frac{\Omega  (N - 1)}{2}}  \frac{\sin\left( \frac{N  \Omega}{2} \right)}{\sin\left( \frac{\Omega}{2} \right)} \label{eq:10.60}
\end{align}
Der Betrag des Spektrums \(W(Ω)\) ist für unterschiedliche Beobachtungszeiten in Bild \ref{fig:10.14} dargestellt.

\begin{figure}
  \centering
  \includegraphics[alt={Spektrum der rechteckförmigen Fensterfunktion \(w(t)\) für Beobachtungszeiten von\\
\(T_{B1}\) = 10T\textsubscript{A} und \(T_{B2}\) = 50T\textsubscript{A}}]{new media/12_14_spektrum_der_rechteckfrmigen_f.png}
  \caption{Spektrum der rechteckförmigen Fensterfunktion \(w(t)\) für Beobachtungszeiten von\\
\(T_{B1}\) = 10T\textsubscript{A} und \(T_{B2}\) = 50T\textsubscript{A}}
  \label{fig:10.14}
\end{figure}

Der Effekt der Zeitbegrenzung ist von der Beobachtungszeit \(T_{B}\) abhängig. Die ersten Nulldurchgänge liegen jeweils an den Stellen
\begin{align}
\Omega &= \frac{2  \pi}{N} \label{eq:10.61}
\end{align}
Mit wachsender Anzahl von Abtastwerten N sinkt demnach die Breite des Spektrums. Das Spektrum nähert sich mit wachsender Anzahl N von Abtastwerten einer Impulsfunktion, die bei der Faltungsoperation das Spektrum unverändert lassen würde. Die Auswirkung der Fensterung wird am Beispiel eines harmonischen Signals diskutiert.
\begin{application}{Leakage-Effekt}
  Im Online-Portal \emph{Systemtheorie Online} verdeutlicht die Applikation \emph{Leakage-Effekt} den Zusammenhang zwischen Fensterung im Zeitbereich und Auswirkung der Fensterung auf den Frequenzbereich.
  \tcblower
  \url{h-ka.de/eit/digitale-signalverarbeitung/übungen/kapitel2}
\end{application}
\subparagraph{Beispiel: Diskrete-Fourier-Transformation einer harmonischen Signalfolge}\label{beispiel-diskrete-fourier-transformation-einer-harmonischen-signalfolge}

Der Effekt der Fensterung wird am Beispiel einer harmonischen Signalfolge erläutert. Sie ist für eine Periodendauer \(N_{0} = 10\) definiert als
\begin{align}
x[ k] &= \cos\left( \frac{2  \pi}{N_{0}}  k \right) = \cos\left( \Omega_{0}  k \right) \label{eq:10.62}
\end{align}
Die Signalfolge weist das Spektrum
\begin{align}
X(\Omega) &= \pi  \left( \delta\left( \Omega + \Omega_{0} \right) + \delta\left( \Omega - \Omega_{0} \right) \right) \label{eq:10.63}
\end{align}
auf. Signalfolge und Spektrum sind für \(Ω_{0} = π\)/5 in Bild \ref{fig:10.15} dargestellt.

\begin{figure}
  \centering
  \includegraphics[alt={Darstellung einer zeitdiskreten Kosinusfunktion und ihres Spektrums}]{new media/12_15_darstellung_einer_zeitdiskrete.png}
  \caption{Darstellung einer zeitdiskreten Kosinusfunktion und ihres Spektrums}
  \label{fig:10.15}
\end{figure}

Zur Berechnung des Spektrums mit der Diskreten-Fourier-Transformation muss die Signalfolge begrenzt werden. Dazu wird das Signal mit einer Rechteckfolge \(w[k]\) multipliziert. Bei einer Begrenzung des Abtastvorgangs auf 2 volle Periodendauern (\(N = 20\) = \SI{2}{\newton}\textsubscript{0}) wird das Spektrum des Signals \(X(Ω)\) mit dem Spektrum der zugehörigen Fensterfunktion \(W(Ω)\) gefaltet. Weil das Spektrum des Signals \(X(Ω)\) aus zwei Impulsen besteht, wird das Spektrum der zugehörigen Fensterfunktion \(W(Ω)\) an die Stelle der Impulse gefaltet. Damit ergibt sich Bild \ref{fig:10.16}.

\begin{figure}
  \centering
  \includegraphics[alt={Darstellung einer zeitbegrenzten Kosinus-Folge und ihrer DFT im Basisband (\(N = 20\))}]{new media/12_16_darstellung_einer_zeitbegrenzt.png}
  \caption{Darstellung einer zeitbegrenzten Kosinus-Folge und ihrer DFT im Basisband (\(N = 20\))}
  \label{fig:10.16}
\end{figure}

Das Spektrum der Fensterfunktion hat keine Auswirkungen auf die berechnete Diskrete-Fourier-Transformierte, weil es an allen Stellen \(Ω_{n}\) ≠ \(Ω_{0}\) den Wert 0 aufweist. Weiterhin ist der Wert \(Ω_{0}\) ein Vielfaches der Auflösung ∆Ω, die entsprechende Amplitude wird deshalb bei der Diskreten-Fourier-Transformierten wirklich berechnet. Bild \ref{fig:10.17} zeigt außerdem, dass eine periodische Fortsetzung der Signalfolge zur ursprünglichen nicht zeitbegrenzten Signalfolge führen würde.

\begin{figure}
  \centering
  \includegraphics[alt={Zeitlich begrenzte Signalfolge \(x_{w}\){[}k{]} und die periodische Fortsetzung}]{new media/12_17_zeitlich_begrenzte_signalfolge.png}
  \caption{Zeitlich begrenzte Signalfolge \(x_{w}\){[}k{]} und die periodische Fortsetzung}
  \label{fig:10.17}
\end{figure}

Das ideale Ergebnis der Diskreten-Fourier-Transformation ist darauf zurückzuführen, dass in diesem Beispiel die Anzahl von Abtastwerten ein Vielfaches der Periodendauer ist. Führt die periodische Wiederholung von Signalausschnitten wieder zur ursprünglichen Signalfolge, entspricht die Diskrete-Fourier-Transformierte dem Spektrum der ursprünglichen Signalfolge \(x_{D}\){[}k{]}. Dieser Fall ist im Allgemeinen aber nicht zu erwarten, da das Signal \(x_{D}\){[}k{]} typischerweise kein harmonisches Signal mit einer bekannten Frequenz sein wird. Bild \ref{fig:10.18} zeigt Signalverläufe für den Fall, dass die Beobachtungsdauer kein Vielfaches der Periodendauer ist (\(N = 15\)). Die periodische Fortsetzung des Signals hat eine Sprungstelle und die Frequenzwerte \(Ω_{n}\), an denen die DFT berechnet wird, ändern sich aufgrund der geänderten Anzahl N von Abtastwerten.
\begin{align}
\Omega_{n} &= \frac{2  \pi  n}{N} \label{eq:10.64}
\end{align}
Damit entspricht kein Wert \(Ω_{n}\) dem eigentlich vorliegenden Wert \(Ω_{0}\), die DFT ist deshalb nur eine Näherung des erwarteten Spektrums.

\begin{figure}
  \centering
  \includegraphics[alt={Darstellung einer zeitbegrenzten Kosinus-Folge und ihrer DFT (\(N = 15\))}]{new media/12_18_darstellung_einer_zeitbegrenzt.png}
  \caption{Darstellung einer zeitbegrenzten Kosinus-Folge und ihrer DFT (\(N = 15\))}
  \label{fig:10.18}
\end{figure}

Weil die ursprünglichen scharfen Spektrallinien des harmonischen Signals sich auf die vorhandenen Stützstellen an den Stellen \(Ω_{n}\) verteilen, wird dieser Effekt als Leck-Effekt (Leakage) bezeichnet. Bei der Diskussion des Dirichlet-Kerns wird darauf hingewiesen, dass mit einer längeren Beobachtungsdauer die Breite des Dirichlet-Kerns reduziert wird. Das Spektrum der Diskreten-Fourier-Transformation wird schmaler. Deshalb zeigt Bild \ref{fig:10.19} die Signalverläufe für den Fall einer längeren Beobachtungsdauer, die kein Vielfaches der Periodendauer ist (\(N = 35\)).

\begin{figure}
  \centering
  \includegraphics[alt={Darstellung einer zeitbegrenzten Kosinus-Folge und ihrer DFT (\(N = 35\))}]{new media/12_19_darstellung_einer_zeitbegrenzt.png}
  \caption{Darstellung einer zeitbegrenzten Kosinus-Folge und ihrer DFT (\(N = 35\))}
  \label{fig:10.19}
\end{figure}

Die einzelnen Werte der Diskreten-Fourier-Transformierten liegen auf dem Spektrum \(X_{w}(Ω)\), das wegen der längeren Beobachtungszeit jedoch deutlich schmaler ist. Außerdem ist mit der längeren Beobachtungsdauer die spektrale Auflösung gestiegen, was in den reduzierten Abstand ∆Ω der Stützstellen resultiert.

▪

Der hier an dem Beispiel eines harmonischen Signals dargestellte Leakage-Effekt gilt sinngemäß für beliebige Signale. Die Beispiele in Abschnitt \ref{sec:10.1} sind so gewählt, dass sich bei der periodischen Wiederholung keine Sprungstellen ergeben und sich der Leakage-Effekt praktisch nicht bemerkbar macht.

\subsection{Alternative Fensterfunktion zur Reduzierung des Leakage-Effektes}\label{sec:alternative-fensterfunktion-zur-reduzierung-des-leakage-effektes}

Der Leakage-Effekt kann darauf zurückgeführt werden, dass die Übergänge bei der periodischen Fortsetzung des zu analysierenden Signals nicht glatt sind. Eine alternative Methode den Leakage-Effekt zu verringern liegt darin, die Übergänge der Signalausschnitte künstlich zu glätten. Dazu zeigt Bild \ref{fig:10.20} die periodische Fortsetzung einer Signalsequenz mit Rechteck-Fenster und mit einem sogenannten Hann-Fenster.

\begin{figure}
  \centering
  \includegraphics[alt={Periodische Fortsetzung einer Signalsequenz mit und ohne Fensterung}]{new media/12_20_periodische_fortsetzung_einer.png}
  \caption{Periodische Fortsetzung einer Signalsequenz mit und ohne Fensterung}
  \label{fig:10.20}
\end{figure}

Das zu analysierende Signal wird mit der Hann-Fensterfunktion multipliziert. Dadurch, dass die Fensterfunktion am Beginn und am Ende des Beobachtungszeitraums stetig auf den Wert null zugeht, wird der Übergang zwischen dem Signal und seiner periodische Fortsetzung geglättet.
\begin{align}
x_{W}[ k] &= \ x_{D}[ k]  w[ k] \label{eq:10.65}
\end{align}
Nach den Rechenregeln zur Fourier-Transformation berechnet sich das Spektrum \(X_{W}(Ω)\) des zeitlich begrenzten Signals \(x_{W}\){[}k{]} aus der Faltung der Spektren \(X_{D}(Ω)\) und \(W(Ω)\). Ideale wäre ein Spektrum, das aus einem Impuls an der Stelle \(Ω = 0\) besteht, da die Faltung des Spektrums \(X_{D}(Ω)\) mit dem Impuls \(δ(Ω)\) das Spektrum unverändert ließe. Allerdings entspricht dem Spektrum \(W(Ω)\) = \(δ(Ω)\) im Zeitbereich die Fensterfunktion \(w[k]\) = 1, die unendlich lang ist und damit zu einer unendlich langen Beobachtungszeit führen würde. Es muss also eine Fensterfunktion \(w[k]\) gefunden werden, die zum einen endlich ist und zum anderen ein Spektrum aufweist, das einer Impulsfunktion \(δ(Ω)\) möglichst nahekommt. Zur Diskussion unterschiedlicher Fensterfunktionen ist in Bild \ref{fig:10.21} das Spektrum des Rechteck-Fensters dargestellt.

\begin{figure}
  \centering
  \includegraphics[alt={Kriterien zur Bewertung von Spektrum von Fensterfunktionen}]{new media/12_21_kriterien_zur_bewertung_von_sp.png}
  \caption{Kriterien zur Bewertung von Spektrum von Fensterfunktionen}
  \label{fig:10.21}
\end{figure}

Das Spektrum besteht aus einem Hauptschwinger (Main-Lobe) und den Nebenschwingern (Side-Lobes). Die Spektren realer Fensterfunktionen haben eine endliche Breite des Hauptmaximums, und es existieren Nebenmaxima, die endlich hohe Amplituden aufweisen. Zur Charakterisierung von Fensterfunktionen werden folgende Bewertungskriterien verwendet.

\begin{itemize}
\item
  Breite des Hauptmaximums\\
  Die Breite des Hauptmaximums ist der Frequenzbereich zwischen \(Ω = 0\) und dem ersten Nulldurchgang des Spektrums \(Ω_{S}\). Eine gute Fensterfunktion hat eine geringe Breite des Hauptmaximums, da dadurch bei der Faltung das Spektrum der gefensterten Funktion weniger stark verbreitert wird.
\item
  Relative Amplitude des ersten Nebenmaximums\\
  Die relative Amplitude des Nebenmaximums ist das Verhältnis von Maximum des Nebenmaximums \(A_{SL}\) zu dem Hauptmaximum \(A_{ML}\). Die relative Amplitude wird in dB angegeben.
\end{itemize}
\begin{align}
a_{REL} &= 20  \log\left( \frac{A_{SL}}{A_{ML}} \right) \label{eq:10.66}
\end{align}
Zielwert ist eine geringe Höhe des Nebenmaximums, also eine kleine relative Amplitude des Nebenmaximums.

Mit der Bestimmung des idealen Fensters beschäftigen sich viele Veröffentlichungen. Es zeigt sich, dass ein ideales Fenster nicht existiert, sondern dass ein Kompromiss zwischen endlicher Breite und relativer Amplitude des Nebenmaximums eingegangen werden muss. Einige Fenster werden im Folgenden beschrieben. Dabei wird zur Vereinfachung der Darstellung von einer Fensterfunktion ausgegangen, die symmetrisch zum Zeitpunkt \(k = 0\) liegt. Diese Fenster sind nicht kausal, können aber durch eine Zeitverschiebung in kausale Fenster überführt werden.

\subsubsection{Rechteckfenster}\label{sec:rechteckfenster-1}

Das Rechteckfenster wird bei gerader Filterordnung N im Zeitbereich über die Gleichung
\begin{align}
w_{REC}[ k] &= \sigma\left( k + \frac{N - 1}{2} \right) - \sigma\left( k - \frac{N + 1}{2} \right) \label{eq:10.67}
\end{align}
beschrieben und hat als Spektrum den sogenannten Dirichlet-Kern
\begin{align}
W_{REC}(\Omega) &= \frac{\sin\left( \frac{N  \Omega}{2} \right)}{\sin\left( \frac{\Omega}{2} \right)} \label{eq:10.68}
\end{align}
Das Rechteckfenster und der normierte Amplitudengang sind in Bild \ref{fig:10.22} dargestellt.

\begin{figure}
  \centering
  \includegraphics[alt={Rechteckfenster und Spektrum des Rechteckfensters für \(N = 20\)}]{new media/12_22_rechteckfenster_und_spektrum_d.png}
  \caption{Rechteckfenster und Spektrum des Rechteckfensters für \(N = 20\)}
  \label{fig:10.22}
\end{figure}

Das Hauptmaximum des Rechteckfensters hat eine Breite von \(Ω_{S} = 2\)·π/N. Die relative Amplitude des Nebenmaximums beträgt \(a_{REL}\) = - \SI{13}{\decibel}.

\subsubsection{Dreieck- oder Bartlett-Fenster}\label{sec:dreieck--oder-bartlett-fenster-1}

Das Dreieckfenster wird im Zeitbereich über die Gleichung
\begin{align}
w_{TRI}[ k] &= \left( k + \frac{N - 1}{2} \right)  \sigma\left( k + \frac{N - 1}{2} \right) - 2  k  \sigma[ k] + \left( k - \frac{N - 1}{2} \right)  \sigma\left( k - \frac{N - 1}{2} \right) \label{eq:10.69}
\end{align}
beschrieben und hat das Spektrum
\begin{align}
W_{TRI}(\Omega) &= \frac{\sin^{2}\left( \frac{N  \Omega}{4} \right)}{\sin^{2}\left( \frac{\Omega}{2} \right)} \label{eq:10.70}
\end{align}
Das Dreieckfenster und der normierte Amplitudengang sind in Bild \ref{fig:10.23} dargestellt.

\begin{figure}
  \centering
  \includegraphics[alt={Dreieckfenster und Spektrum des Dreieckfensters für \(N = 20\)}]{new media/12_23_dreieckfenster_und_spektrum_de.png}
  \caption{Dreieckfenster und Spektrum des Dreieckfensters für \(N = 20\)}
  \label{fig:10.23}
\end{figure}

Das Hauptmaximum des Dreieckfensters hat eine Breite von \(Ω_{S} = 4\)·π/N. Die relative Amplitude des Nebenmaximums beträgt \(a_{REL}\) = - \SI{25}{\decibel}.

\subsubsection{Hann-Fenster}\label{sec:hann-fenster-1}

Auch das Hann-Fenster weist am Beginn und am Ende der Beobachtungszeit keine Sprünge auf, sondern geht an diesen Punkten gegen den Wert 0. Der stetige Übergang wird mit einer Kosinusfunktion beschrieben.
\begin{align}
w_{HAN}[ k] &= \left( 0.5 + 0.5  \cos\left( \frac{2  \pi  k}{N} \right) \right)  \left( \sigma\left( k + \frac{N - 1}{2} \right) - \sigma\left( k - \frac{N + 1}{2} \right) \right) \label{eq:10.71}
\end{align}
Die Fensterfunktion ergibt sich aus dem Produkt von Rechteckfenster und der Kosinusfunktion. Das Spektrum des Hann-Fensters kann über die Faltungsoperation berechnet werden. Es ergibt sich zu
\begin{align}
W_{HAN}(\Omega) &= \frac{1}{4}  \frac{\sin\left( \frac{N  \left( \Omega + \frac{2  \pi}{N} \right)}{2} \right)}{\sin\left( \frac{\Omega + \frac{2  \pi}{N}}{2} \right)} + \frac{1}{2}  \frac{\sin\left( \frac{N  \Omega}{2} \right)}{\sin\left( \frac{\Omega}{2} \right)} + \frac{1}{4}  \frac{\sin\left( \frac{N  \left( \Omega - \frac{2  \pi}{N} \right)}{2} \right)}{\sin\left( \frac{\Omega - \frac{2  \pi}{N}}{2} \right)} \label{eq:10.72}
\end{align}
Das Hann-Fenster und sein Spektrum sind in Bild \ref{fig:10.24} dargestellt.

\begin{figure}
  \centering
  \includegraphics[alt={Hann-Fenster und Spektrum des Hann-Fensters für \(N = 20\)}]{new media/12_24_hannfenster_und_spektrum_des_h.png}
  \caption{Hann-Fenster und Spektrum des Hann-Fensters für \(N = 20\)}
  \label{fig:10.24}
\end{figure}

Das Hauptmaximum des Hann-Fensters hat eine Breite von \(Ω_{S} = 4\)·π/N und ist damit doppelt so breit wie bei dem Rechteckfenster. Die relative Amplitude des Nebenmaximums ist mit \(a_{REL}\) = - \SI{32}{\decibel} deutlich größer als bei dem Rechteckfenster.

\subsubsection{Hamming-Fenster}\label{sec:hamming-fenster-1}

Hamming- und Hann-Fenster unterscheiden sich nur geringfügig in der Definition.
\begin{align}
w_{HAM}[ k] &= \left( 0.54 + 0.46  \cos\left( \frac{2  \pi  k}{N} \right) \right)  \left( \sigma\left( k + \frac{N - 1}{2} \right) - \sigma\left( k - \frac{N + 1}{2} \right) \right) \label{eq:10.73}
\end{align}
Auch die Berechnung des Spektrums der beiden Fensterfunktionen ist ähnlich. Für das Hamming-Fenster ergibt sich das Spektrum
\begin{align}
W_{HAM}(\Omega) &= 0.23  \frac{\sin\left( \frac{N  \left( \Omega + \frac{2  \pi}{N} \right)}{2} \right)}{\sin\left( \frac{\Omega + \frac{2  \pi}{N}}{2} \right)} + 0.54  \frac{\sin\left( \frac{N  \Omega}{2} \right)}{\sin\left( \frac{\Omega}{2} \right)} + 0.23  \frac{\sin\left( \frac{N  \left( \Omega - \frac{2  \pi}{N} \right)}{2} \right)}{\sin\left( \frac{\Omega - \frac{2  \pi}{N}}{2} \right)} \label{eq:10.74}
\end{align}
Das Hamming-Fenster und sein Spektrum sind in Bild \ref{fig:10.25} dargestellt.

\begin{figure}
  \centering
  \includegraphics[alt={Hamming-Fenster und Spektrum des Hamming-Fensters für \(N = 20\)}]{new media/12_25_hammingfenster_und_spektrum_de.png}
  \caption{Hamming-Fenster und Spektrum des Hamming-Fensters für \(N = 20\)}
  \label{fig:10.25}
\end{figure}

Das Hauptmaximum des Hamming-Fensters hat eine Breite von \(Ω_{S} = 4\)·π/N. Die relative Amplitude des Nebenmaximums beträgt mit \(a_{REL}\) = - \SI{42}{\decibel} und ist nochmals deutlich größer als bei dem Hann-Fenster.

\subsubsection{Vergleich der Fensterfunktionen}\label{sec:vergleich-der-fensterfunktionen-1}

Tabelle \ref{tab:10.6} stellt die berechneten Kennwerte der unterschiedlichen Fenster zusammen.

\begin{table}
  \centering
  \caption{Tabellarischer Vergleich der Kennwerte unterschiedlicher Fensterfunktionen}
  \label{tab:10.6}
  \begin{tabular}{lll}
    \toprule
    Fenster- 
Funktion
 & 
Breite des 
Hauptmaximums \(Ω_{S}\)
 & 
Relative Amplitude 
der Nebenmaxima \(a_{REL}\) \\
    \midrule
    Rechteck & 2·π / N & - \SI{13}{\decibel} \\
    Dreieck & 4·π / N & - \SI{25}{\decibel} \\
    Hann & 4·π / N & - \SI{32}{\decibel} \\
    Hamming & 4·π / N & - \SI{42}{\decibel} \\
    \bottomrule
  \end{tabular}
\end{table}

Der Vergleich zeigt, dass es eine ideale Fensterfunktion nicht gibt. Das Spektrum des Rechteckfensters ist schmal, aber es weist vergleichsweise hohe Nebenmaxima auf. Demgegenüber haben die Spektren von Hann- und Hamming-Fenster relativ kleine Nebenmaxima, die Hauptmaxima sind aber breiter als bei dem Rechteckfenster. Für praktische Anwendungen bietet das Hamming-Fenster einen guten Kompromiss zwischen Breite des Hauptmaximums und relativer Amplitude des Nebenmaximums.

\begin{example}{Vergleich verschiedener Fensterfunktionen bei der Diskreten-Fourier-Transformation}
\label{beispiel-vergleich-verschiedener-fensterfunktionen-bei-der-diskreten-fourier-transformation}

Als Beispiel wird eine harmonische Funktion mit der normierten Kreisfrequenz \(Ω_{0} = 2\)·π/7 abgetastet. Es werden \(N = 32\) Abtastwerte aufgenommen, wobei die Signale mit Rechteck-, Hann- und Hamming-Fenster multipliziert werden. Bild \ref{fig:10.26} zeigt die aufgenommenen Werte.

\begin{fixedfigure}
  \includegraphics[alt={Abtastwerte der harmonischen Funktion nach Gewichtung mit Rechteck-, Hann-\\
und Hamming-Fenster (Frequenz \(Ω_{0} = 2\)·π/7)}]{new media/12_26_abtastwerte_der_harmonischen_f.png}
  \captionof{figure}{Abtastwerte der harmonischen Funktion nach Gewichtung mit Rechteck-, Hann-\\
und Hamming-Fenster (Frequenz \(Ω_{0} = 2\)·π/7)}
  \label{fig:10.26}
\end{fixedfigure}

Es ist deutlich zu erkennen, dass die Werte an den Rändern des Beobachtungszeitraums auf null reduziert werden. Mit den jeweils 32 Werten wird eine Diskrete-Fourier-Transformation durchgeführt. Bild \ref{fig:10.27} zeigt die Ergebnisse.

\begin{fixedfigure}
  \includegraphics[alt={Diskrete-Fourier-Transformierte der harmonischen Funktion\\
nach Gewichtung mit Rechteck- und Hann- und Hamming-Fenster (Frequenz \(Ω_{0} = 2\)·π/7)}]{new media/12_27_diskretefouriertransformierte.png}
  \captionof{figure}{Diskrete-Fourier-Transformierte der harmonischen Funktion\\
nach Gewichtung mit Rechteck- und Hann- und Hamming-Fenster (Frequenz \(Ω_{0} = 2\)·π/7)}
  \label{fig:10.27}
\end{fixedfigure}

Kein Fenster führt zur Berechnung des Spektrums, das sich aus der analytischen Rechnung ergibt. Da an der entsprechenden Frequenz \(Ω_{0}\) keine Stützstelle liegt, ist dies auch nicht möglich. Das Spektrum ist bei allen Fenstern ähnlich auf die beiden Stützstellen verteilt, die in direkter Nachbarschaft zum Impuls an der Frequenz \(Ω_{0}\) liegen.

Bei dem Rechteck-Fenster fällt der Betrag des Spektrums nur langsam ab. Stützstellen, die sehr weit weg von der Frequenz \(Ω_{0}\) liegen, weisen noch einen erheblichen Amplitudenwert auf. Bei Hann- und Hamming-Fenster fallen die Beträge des Spektrums deutlich schneller ab. Das Beispiel bestätigt damit die theoretische Betrachtung in diesem Abschnitt.

Zum Vergleich zeigt Bild \ref{fig:10.28} die über die Diskrete-Fourier-Transformation berechneten Spektren bei einer Frequenz \(Ω_{0} = 2\)·π/8.

\begin{fixedfigure}
  \includegraphics[alt={Diskrete-Fourier-Transformierte der harmonischen Funktion\\
nach Gewichtung mit Rechteck- und Hann- und Hamming-Fenster (Frequenz \(Ω_{0} = 2\)·π/8)}]{new media/12_28_diskretefouriertransformierte.png}
  \captionof{figure}{Diskrete-Fourier-Transformierte der harmonischen Funktion\\
nach Gewichtung mit Rechteck- und Hann- und Hamming-Fenster (Frequenz \(Ω_{0} = 2\)·π/8)}
  \label{fig:10.28}
\end{fixedfigure}

In diesem Fall wird das ideale Spektrum mit dem Rechteckfenster richtig identifiziert. Die beiden übrigen Fenster verbreitern das Spektrum aufgrund der größeren Breite des Hauptmaximums. Der Einsatz der Fensterfunktionen verbessert das Ergebnis offensichtlich nicht in allen Fällen.

\end{example}

Zusammenfassend können damit zur Reduktion des Leakage-Effektes folgende Maßnahmen ergriffen werden:

\begin{itemize}
\item
  Anpassung der Abtast- und Beobachtungszeit an die Aufgabenstellung\\
  Im Idealfall ist die Beobachtungszeit so gewählt, dass ein von sich aus periodisch fortgesetztes Signal entsteht. Die Maßnahme ist sehr wirkungsvoll, lässt sich aber in den wenigsten Anwendungen umsetzen, da die zu analysierenden Signale im Allgemeinen unbekannt sind.
\item
  Vergrößerung des Beobachtungszeitraums\\
  Eine Vergrößerung des Beobachtungsintervalls führt zu einer Verringerung der Breite des Hauptschwingers, wodurch die spektrale Auflösung der Diskreten-Fourier-Transformation ∆Ω verbessert wird.
\item
  Anwendung von Fensterfunktionen\\
  Durch Fensterfunktionen werden die bei dem Rechteckfenster vergleichsweise schwach gedämpften Nebenschwinger stärker gedämpft. Dieser Vorteil wird jedoch durch eine Verdoppelung der Breite des Hauptschwingers erkauft.
\end{itemize}

\section{Frequenzauflösung der Diskreten-Fourier-Transformation -- Zero-Padding}\label{sec:10.4}

Das Einfügen von Nullen in die Signalfolge oder in das über die Diskrete-Fourier-Transformation berechnete Spektrum führt zur Interpolation des kontinuierlichen Spektrums beziehungsweise der entsprechenden kontinuierlichen Zeitfunktion. Beide Verfahren werden im Folgenden diskutiert.

\subsection{Zero-Padding im Zeitbereich}\label{sec:zero-padding-im-zeitbereich}

In Abschnitt \ref{sec:10.1} wird gezeigt, dass die Auflösung der Diskreten-Fourier-Transformation von der Länge des Beobachtungsintervalls \(T_{B}\) beziehungsweise der Anzahl von Abtastwerten N abhängt.
\begin{align}
\Delta\omega_{N} &= \frac{2  \pi}{T_{A}  N} = \frac{2  \pi}{T_{B}} \label{eq:10.75}
\end{align}
Liegen keine zusätzlichen Abtastwerte vor, können die vorhandenen Messwerte mit M Nullen ergänzt werden. Es entsteht eine Signalfolge der Länge N + M. Aufgrund der größeren Anzahl von Folgenwerten verbessert sich die numerische Auflösung zu
\begin{align}
\Delta\omega_{N + M} &= \frac{2  \pi}{T_{A}  (N + M)} < \frac{2  \pi}{T_{A}  N} = \Delta\omega_{N} \label{eq:10.76}
\end{align}
Dieser Vorgang wird als Zero-Padding bezeichnet. Um die Auswirkung des Zero-Padding auf die Diskrete-Fourier-Transformierte bewerten zu können, werden die Fourier-Transformierten unterschiedlicher Signalfolgen betrachtet. Es werden drei Fälle diskutiert.

\begin{itemize}
\item
  Fall 1: Standard DFT mit \(N_{1} = 16\) Abtastwerten\\
  Im ersten Fall wird die Signalfolge mit einem Fenster der Länge \(N_{1} = 16\) beobachtet und die Diskrete-Fourier-Transformierte berechnet.
\item
  Fall 2: Zero-Padding\\
  Im zweiten Fall wird die Länge der Signalfolge auf \(N_{1}\) + \(M = 32\) verdoppelt, indem die gefensterte Signalfolge mit \(M = 16\) Nullen aufgefüllt wird.
\item
  Fall 3: DFT mit \(N_{2} = 32\) Abtastwerten\\
  Im dritten Fall wird die Beobachtungszeit auf \(N_{2} = 32\) Abtastwerte erhöht.
\end{itemize}

Bild \ref{fig:10.29} zeigt die Berechnungsergebnisse aller drei Fälle. Ein Vergleich der beiden ersten Fälle zeigt, dass durch die Erweiterung der Signalfolge mit Nullen das Spektrum erwartungsgemäß enger abgetastet wird. Da die Fensterung des Signals jedoch nicht verändert wird, ändert sich die Breite des Dirichlet-Kerns nicht. Aus diesem Grund wird das Spektrum \(X_{W1}(Ω)\) nur mit einer höheren Auflösung abgetastet. Zero-Padding eignet sich damit zur Interpolation des Spektrums, es eignet sich nicht dazu, den durch die Fensterung entstehenden Leakage-Effekt zu reduzieren. Das Einfügen von Nullen führt damit nicht zu einem Informationsgewinn. Zero-Padding kann dazu eingesetzt werden, um das Spektrum besser darzustellen und um lokale Maxima genauer zu identifizieren.

Im Gegensatz dazu zeigt der Vergleich der Fälle 1 und 3, dass durch eine Vergrößerung des Beobachtungsintervalls die Breite des Dirichlet-Kerns verringert wird, was zu einer Reduzierung des Leakage-Effektes führt. Der Abstand der Stützstellen ∆Ω entspricht in Fall 3 dem Fall 2.

\begin{figure}
  \centering
  \includegraphics[alt={Vergleich der Diskreten-Fourier-Transformierten\\
Fall 1: Signalfolge der Länge \(N_{1} = 16\)\\
Fall 2: Signalfolge der Länge \(N_{1}\) mit Zero-Padding \(M = 16\)\\
Fall 3: Signalfolge der Länge \(N_{2}\) = \SI{2}{\newton}\textsubscript{1} = 32}]{new media/12_29_vergleich_der_diskretenfourier.png}
  \caption{Vergleich der Diskreten-Fourier-Transformierten\\
Fall 1: Signalfolge der Länge \(N_{1} = 16\)\\
Fall 2: Signalfolge der Länge \(N_{1}\) mit Zero-Padding \(M = 16\)\\
Fall 3: Signalfolge der Länge \(N_{2}\) = \SI{2}{\newton}\textsubscript{1} = 32}
  \label{fig:10.29}
\end{figure}

\subsection{Rekonstruktion der Fourier-Transformierten für Signalfolgen aus der DFT}\label{sec:rekonstruktion-der-fourier-transformierten-fuxfcr-signalfolgen-aus-der-dft}

Zur Berechnung der Diskreten-Fourier-Transformation wird eine Signalfolge \(x_{W}\){[}k{]} verwendet, die ein abgetastetes Signal in einem Beobachtungsfenster darstellt. Durch die Berechnung der Diskreten-Fourier-Transformierten wird die Signalfolge periodisch fortgesetzt. Das Spektrum der periodisch fortgesetzten Signalfolge \(x_{P}\){[}k{]} ist die Diskrete-Fourier-Transformierte \(X_{n}\). Sie ist nur an den Punkten \(Ω_{n}\) definiert und kann über Impulsfunktionen dargestellt werden als
\begin{align}
X_{P}(\Omega) &= \sum\limits_{n = 0}^{N - 1}{X_{n}  \delta\left( \Omega - \frac{2  \pi  n}{N} \right)} \label{eq:10.77}
\end{align}
Um das Spektrum der gefensterten Signalfolge \(x_{W}\){[}k{]} ohne periodische Wiederholung zu berechnen, muss von der Beschreibung im Zeitbereich auf den Frequenzbereich geschlossen werden. Im Zeitbereich gilt der Zusammenhang
\begin{align}
x_{W}[ k] &= x_{P}[ k]  w[ k] = x_{P}[ k]  \left( \sigma[ k] - \sigma[ k - N] \right) \label{eq:10.78}
\end{align}
Der Multiplikation im Zeitbereich entspricht die Faltung im Frequenzbereich. Das Spektrum der Fensterfunktion \(W(Ω)\) lautet
\begin{align}
W(\Omega) &= e^{- j  \Omega  \frac{N - 1}{2}}  \frac{\sin\left( \frac{\Omega  N}{2} \right)}{\sin\left( \frac{\Omega}{2} \right)} \label{eq:10.79}
\end{align}
Da das Spektrum der periodisch fortgesetzten Folge ausschließlich aus gewichteten Impulsen besteht, wird durch die Faltung das Spektrum an die Stellen verschoben, an denen die Impulse liegen. Es ergibt sich
\begin{align}
X_{W}(\Omega) &= X_{P}(\Omega)*W(\Omega) = \sum\limits_{n = 0}^{N - 1}{X_{n}  e^{- j  \left( \Omega - \frac{2  \pi  n}{N} \right)  \frac{N - 1}{2}}  \frac{\sin\left( \frac{\left( \Omega - \frac{2  \pi  n}{N} \right)  N}{2} \right)}{\sin\left( \frac{\Omega - \frac{2  \pi  n}{N}}{2} \right)}} \label{eq:10.80}
\end{align}
Das Spektrum des gefensterten Signals \(x_{W}\){[}k{]} besitzt das kontinuierliche Spektrum \(X_{W}(Ω)\).

\subsection{Zero-Padding im Frequenzbereich}\label{sec:zero-padding-im-frequenzbereich}

\textless\textless\textless{} wird später ergänzt \textgreater\textgreater\textgreater{}

\section{Fast-Fourier-Transformation}\label{sec:10.5}

Die Berechnung der Diskreten-Fourier-Transformation ist Basis von Filteralgorithmen und der Spektralschätzung. Aus diesem Grund sind Algorithmen zur Berechnung der Diskreten-Fourier-Transformation notwendig, die hinsichtlich des Realisierungsaufwands optimiert sind. Sie werden als Fast-Fourier-Transformation (FFT) bezeichnet. Mit Hilfe der FFT-Algorithmen wird die Anzahl von Rechenoperationen von N² auf Nlog\textsubscript{2}(N) reduziert. Zum Beispiel wird durch Einsatz einer Fast-Fourier-Transformation für \(N = 1024\) die Anzahl von Rechenoperationen gegenüber der Diskreten-Fourier-Transformation auf 1 \% reduziert. Dieser Abschnitt stellt die Grundidee vor und beschreibt eine Implementierungsform.

\subsection{Matrizendarstellung der Diskreten-Fourier-Transformation}\label{sec:matrizendarstellung-der-diskreten-fourier-transformation}

Ausgangspunkt für die Fast-Fourier-Transformation ist die Matrizendarstellung der Diskreten-Fourier-Transformation. Dabei wird die Definitionsgleichung
\begin{align}
X_{n} &= \sum\limits_{k = 0}^{N - 1}{x[ k]  e^{- j  \left( 2  \pi  \frac{n}{N}  k \right)}}\ \label{eq:10.81}
\end{align}
in eine Matrizenschreibweise überführt. Die einzelnen Werte der Diskreten-Fourier-Transformation werden als Spaltenvektor \ul{X} geschrieben. Er besitzt N Elemente.
\begin{align}
\overset{\underbar{}}{\mathbf{X}} &= \begin{pmatrix}
X_{0} \\
X_{1} \\
 \vdots \\
X_{N - 1}
\end{pmatrix} \label{eq:10.82}
\end{align}
Auch die Abtastwerte \(x[k]\) können als Spaltenvektor \ul{x} mit N Elementen dargestellt werden.
\begin{align}
\overset{\underbar{}}{x} &= \begin{pmatrix}
x[ 0] \\
x[ 1] \\
 \vdots \\
x[ N - 1]
\end{pmatrix} \label{eq:10.83}
\end{align}
Die Summe aus Gleichung \eqref{eq:10.81} kann über das Produkt einer Matrix \textbf{W} mit dem Vektor \ul{x} aufgefasst werden. Es ergibt sich
\begin{align}
\overset{\underbar{}}{X} = \begin{pmatrix}
X_{0} \\
X_{1} \\
 \vdots \\
X_{N - 1}
\end{pmatrix} = \begin{pmatrix}
e^{- j  \left( 2  \pi  \frac{0}{N}  0 \right)} & e^{- j  \left( 2  \pi  \frac{0}{N}  1 \right)} & s & e^{- j  \left( 2  \pi  \frac{0}{N}  (N - 1) \right)} \\
e^{- j  \left( 2  \pi  \frac{1}{N}  0 \right)} & e^{- j  \left( 2  \pi  \frac{1}{N}  1 \right)} & s & e^{- j  \left( 2  \pi  \frac{1}{N}  (N - 1) \right)} \\
 \vdots & \vdots & \ddots & \vdots \\
e^{- j  \left( 2  \pi  \frac{N - 1}{N}  0 \right)} & e^{- j  \left( 2  \pi  \frac{N - 1}{N}  1 \right)} & \ldots & e^{- j  \left( 2  \pi  \frac{N - 1}{N}  (N - 1) \right)}
\end{pmatrix}  \begin{pmatrix}
x[ 0] \\
x[ 1] \\
 \vdots \\
x[ N - 1]
\end{pmatrix} \label{eq:10.84}
\end{align}
In der Matrix treten Exponentialfunktionen auf, die eine Drehung in der komplexen Ebene beschreibt. Er wird als komplexer Drehfaktor bezeichnet und zur Erleichterung der Schreibweise abgekürzt zu
\begin{align}
W_{N} &= e^{- j  \left( 2  \pi  \frac{1}{N} \right)} \label{eq:10.85}
\end{align}
Bild \ref{fig:10.30} stellt die komplexen Drehfaktoren für die Ordnung \(N = 7\) und \(N = 8\) in der komplexen Ebene dar.

\begin{figure}
  \centering
  \includegraphics[alt={Darstellung des komplexen Drehfaktors für die Ordnung \(N = 7\) und \(N = 8\) in der komplexen Ebene}]{new media/12_30_darstellung_des_komplexen_dreh.png}
  \caption{Darstellung des komplexen Drehfaktors für die Ordnung \(N = 7\) und \(N = 8\) in der komplexen Ebene}
  \label{fig:10.30}
\end{figure}

Ist die Ordnung N einer Zweipotenz, weisen die komplexen Drehfaktoren eine Symmetrie auf. Diese Symmetrie ist ein Ansatz zur Reduktion des Rechenaufwandes bei der Fast-Fourier-Transformation. Mit der Definition des komplexen Drehfaktors kann die Matrizendarstellung überführt werden zu
\begin{align}
\overset{\underbar{}}{X} = \begin{pmatrix}
X_{0} \\
X_{1} \\
 \vdots \\
X_{N - 1}
\end{pmatrix} = \begin{pmatrix}
W_{N}^{0} & W_{N}^{0} & s & W_{N}^{0} \\
W_{N}^{0} & W_{N}^{1} & s & W_{N}^{N - 1} \\
 \vdots & \vdots & \ddots & \vdots \\
W_{N}^{0} & W_{N}^{N - 1} & s & W_{N}^{(N - 1)^{2}}
\end{pmatrix}  \begin{pmatrix}
x[ 0] \\
x[ 1] \\
 \vdots \\
x[ N - 1]
\end{pmatrix} = \mathbf{F}  \overset{\underbar{}}{x} \label{eq:10.86}
\end{align}
Die Matrix \textbf{F} wird als Fourier-Matrix bezeichnet. Entsprechend gilt für die inverse Diskrete-Fourier-Transformation
\begin{align}
x[ k] &= \frac{1}{N}  \sum\limits_{n = 0}^{N - 1}X_{n}  e^{j  \left( 2  \pi  \frac{n}{N}  k \right)} \label{eq:10.87}
\end{align}
in Matrix-Schreibweise
\begin{align}
{\underline{x}} &=  {\frac{1}{N}  \left(\begin{array}{cccc} {W_{N}^{0} } & {W_{N}^{0} } & {s } & {W_{N}^{0} } \\
{W_{N}^{0} } & {W_{N}^{-1} } & {s } & {W_{N}^{-\left(N-1\right)} } \\
{\vdots } & {\vdots } & {\ddots } & {\vdots } \\
{W_{N}^{0} } & {W_{N}^{-\left(N-1\right)} } & {s } & {W_{N}^{-\left(N-1\right)^{2} } } \end{array}\right) \left(\begin{array}{c} {X_{0} } \\
{X_{1} } \\
{\vdots } \\
{X_{N-1} } \end{array}\right)} \\
{} &=  {\left(\begin{array}{cccc} {W_{N}^{0} } & {W_{N}^{0} } & {s } & {W_{N}^{0} } \\
{W_{N}^{0} } & {W_{N}^{1} } & {s } & {W_{N}^{N-1} } \\
{\vdots } & {\vdots } & {\ddots } & {\vdots } \\
{W_{N}^{0} } & {W_{N}^{N-1} } & {s } & {W_{N}^{\left(N-1\right)^{2} } } \end{array}\right)^{-1}  \left(\begin{array}{c} {X_{0} } \\
{X_{1} } \\
{\vdots } \\
{X_{N-1} } \end{array}\right)=\boldsymbol{\mathrm{F}}^{-1}  \underline{X}}
 \label{eq:10.88}
\end{align}
\subsection{Grundidee der Fast-Fourier-Transformation}\label{sec:grundidee-der-fast-fourier-transformation}

Die Matrizenschreibweise zeigt, dass die Diskrete-Fourier-Transformation der Länge N mit NN komplexe Multiplikationen und N(N - 1) komplexen Additionen berechnet wird. Ein Kernpunkt der Aufwandsreduktion ist, die Periodizität der komplexen Exponentialfunktion beziehungsweise des komplexen Drehfaktors \(W_{N}\) auszunutzen. Unter Berücksichtigung der Periodizität
\begin{align}
W_{N}^{n} &= W_{N}^{n + N} \label{eq:10.89}
\end{align}
kann im Fall einer Ordnung \(N = 4\) die Fourier-Matrix geschrieben werden als
\begin{align}
\begin{pmatrix}
W_{4}^{0} & W_{4}^{0} & W_{4}^{0} & W_{4}^{0} \\
W_{4}^{0} & W_{4}^{1} & W_{4}^{2} & W_{4}^{3} \\
W_{4}^{0} & W_{4}^{2} & W_{4}^{4} & W_{4}^{6} \\
W_{4}^{0} & W_{4}^{3} & W_{4}^{6} & W_{4}^{9}
\end{pmatrix} = \begin{pmatrix}
W_{4}^{0} & W_{4}^{0}  W_{4}^{0} & W_{4}^{0} & W_{4}^{0}  W_{4}^{0} \\
W_{4}^{0} & W_{4}^{1}  W_{4}^{0} & W_{4}^{2} & W_{4}^{1}  W_{4}^{2} \\
W_{4}^{0} & W_{4}^{2}  W_{4}^{0} & W_{4}^{0} & W_{4}^{2}  W_{4}^{0} \\
W_{4}^{0} & W_{4}^{3}  W_{4}^{0} & W_{4}^{2} & W_{4}^{3}  W_{4}^{2}
\end{pmatrix} \label{eq:10.90}
\end{align}
Ziel ist es, die in rot eingezeichneten gemeinsamen Vektoren für eine Zusammenfassung zu nutzen. Das kann dadurch geschehen, dass die entsprechenden Variablen vorab addiert und erst dann mit dem komplexen Drehfaktor multipliziert werden. Der resultierende Signalfluss ist in Bild \ref{fig:10.31} dargestellt.

\begin{figure}
  \centering
  \includegraphics[alt={Rechenschema der FFT für die Ordnung \(N = 4\)}]{new media/12_31_rechenschema_der_fft_fr_die_or.pdf}
  \caption{Rechenschema der FFT für die Ordnung \(N = 4\)}
  \label{fig:10.31}
\end{figure}

In der ersten Stufe werden die gewichteten Summen von Spalte 1 und 3 beziehungsweise 2 und 4 gebildet. In der zweiten Stufe wird die Multiplikation mit dem komplexen Drehfaktor vorgenommen. Es werden insgesamt 8 Multiplikationen und 8 Additionen benötigt. Mit diesem Vorgehen lässt sich sowohl die Anzahl der komplexen Additionen, als auch Multiplikationen von \(N^{2}\) auf Nlog\textsubscript{2}(N) reduzieren. Diese Symmetrie ist allerdings nur vorhanden, wenn die Anzahl von Messwerten eine Potenz von 2 ist.
\begin{align}
N &= 2^{n} \label{eq:10.91}
\end{align}
Kann diese Bedingung nicht erreicht werden, werden zu den vorhandenen Messwerten solange Nullen hinzugefügt, bis die entstehende Signalfolge eine Zweierpotenz aufweist. Eine weitere Aufwandsreduzierung lässt sich im Fall reeller Zahlenfolgen erreichen. Hierzu werden die in Abschnitt \ref{sec:10.2.2} dargestellten Symmetrieeigenschaften der Diskreten-Fourier-Transformation genutzt, um statt einer komplexen Signalfolge parallel zwei reelle Zahlenfolgen zu berechnen.

\subsection{Radix-2-FFT-Algorithmus}\label{sec:radix-2-fft-algorithmus}

\textless\textless\textless{} wird später ergänzt \textgreater\textgreater\textgreater{}

\subsection{FFT-Algorithmus für reelle Signale}\label{sec:fft-algorithmus-fuxfcr-reelle-signale}

\section{Anwendungen der Fast-Fourier-Transformation}\label{sec:10.6}

\subsection{Faltungsoperation und ihre Implementierung im Frequenzbereich}\label{sec:faltungsoperation-und-ihre-implementierung-im-frequenzbereich}

\subsection{Blockweise Verarbeitung mit dem Overlap-Add- und Overlapp-Save}\label{sec:blockweise-verarbeitung-mit-dem-overlap-add--und-overlapp-save}

\subsection{Implementierungsvarianten}\label{sec:implementierungsvarianten}
